\documentclass[11pt,a4paper]{article}
\usepackage[T1]{fontenc}
\usepackage[utf8]{inputenc}
\usepackage{lmodern}
\usepackage[a4paper,margin=25mm,headheight=15pt]{geometry}
\usepackage{amsmath,amssymb,amsthm,mathtools}
\usepackage{mathrsfs}% added for Part II (batch 86A)
\usepackage{booktabs,array,longtable}
\usepackage{tabularx}% added for Part II (batch 86A)
\usepackage{microtype}
\usepackage[dvipsnames]{xcolor}
\usepackage{xurl}
\usepackage{enumitem}
\usepackage{fancyhdr}
\usepackage{listings}% added for Part II (batch 86A)
\usepackage{hyperref}
\hypersetup{colorlinks=true,linkcolor=MidnightBlue,citecolor=MidnightBlue,urlcolor=MidnightBlue,
 pdftitle={Proving Two Arithmetic Conjectures for OEIS A321941},
 pdfauthor={ChatGPT},pdfsubject={Integrality, congruences, formal series, asymptotic inversion}}
\pagestyle{fancy}
\fancyhf{}
\fancyhead[L]{\small Arithmetic of OEIS A321941}
\fancyhead[R]{\small Research draft | 1 October 2026}
\fancyfoot[C]{\thepage}
\renewcommand{\headrulewidth}{0.3pt}
\setlength{\parskip}{4pt}
\setlength{\parindent}{0pt}
\setlength{\emergencystretch}{2em}
\setlist[enumerate]{leftmargin=*,itemsep=5pt,topsep=4pt}
\numberwithin{equation}{section}
\newtheorem{theorem}{Theorem}[section]
\newtheorem{lemma}[theorem]{Lemma}
\newtheorem{proposition}[theorem]{Proposition}
\newtheorem{corollary}[theorem]{Corollary}
\theoremstyle{definition}
\newtheorem{definition}[theorem]{Definition}
\theoremstyle{remark}
\newtheorem{remark}[theorem]{Remark}
% Added for Part II (batch 86A): its research questions and listings.
\theoremstyle{definition}
\newtheorem{question}[theorem]{Research question}
\lstset{basicstyle=\ttfamily\scriptsize,columns=fullflexible,
 keepspaces=true,showstringspaces=false,breaklines=true,
 frame=single,rulecolor=\color{black!20},framesep=5pt,
 aboveskip=10pt,belowskip=10pt}
\newcommand{\Z}{\mathbb Z}
\newcommand{\Q}{\mathbb Q}
\newcommand{\R}{\mathbb R}
\newcommand{\Eop}{\mathcal E}
\newcommand{\Oop}{\mathcal O}
\newcommand{\Hop}{\mathcal H}
\newcommand{\coeff}[1]{\left[#1\right]}
\newcommand{\ord}{\operatorname{ord}}
\newcommand{\den}{\operatorname{den}}
\newcommand{\Cat}{\operatorname{Cat}}
\newcommand{\ind}{\mathbf 1}
\newcommand{\dif}{\,\mathrm d}
\newcommand{\repoSHA}{\texttt{0a6d798c5f39775b2fb4d85fef643fef996c87cf}}

\begin{document}
\hypersetup{pageanchor=false}
\begin{titlepage}
\vspace*{9mm}
{\small\sffamily OEIS RESEARCH \quad / \quad SPECIAL FUNCTIONS AND FORMAL ARITHMETIC\par}
\vspace{12mm}
{\LARGE\bfseries Proving Two Arithmetic Conjectures\par}
\vspace{3mm}
{\LARGE\bfseries for OEIS A321941\par}
\vspace{8mm}
{\large A nonlinear product identity, integral polynomial deformations,\par
sharp denominator bounds, and asymptotic inversion\par}
\vspace{10mm}
{\normalsize ChatGPT\par Research draft prepared for the ProveIt research program\par
1 October 2026\par}
\vspace{9mm}
\begin{abstract}
Brent, Glasser, and Guttmann introduced rational coefficients $d_k$ in the
all-orders asymptotic expansion of a Hadamard product involving the exponential
integral. They conjectured that $r_k=64^kd_k$ is always an integer, and observed
that $r_k\equiv\binom{2k}{k}\pmod{32}$ for $k\ge3$. This article proves both
arithmetic assertions. The key is to retain a quadratic invariant of the
underlying second-order recurrence instead of using only its linear
symmetric-square recurrence. After scaling, the invariant gives a triangular
nonlinear equation with integer shift operators. An evenness induction removes
all coefficient-index divisions.

The construction extends to polynomials $P_k(t)\in\Z[t]$ with $P_k(1)=r_k$,
$P_k(t)\in2\Z[t]$ for $k>0$, and a coefficientwise congruence modulo $32$.
Both the leading coefficient and the constant term of $P_k$ have closed forms.
Consequences include the minimality of the scale $64$, sharp denominator bounds
at infinitely many indices, and exact denominator formulas whenever the index
has at most four nonzero binary digits. Finally, the product expansion is
inverted to all algebraic orders; after normalization, the inverse coefficients
have no denominator primes other than $3$.

The analytic existence of the original expansion is the published theorem of
Brent--Glasser--Guttmann. The new arithmetic arguments are proved here. Their
separate negativity conjecture is not proved. Exact finite checks accompany the
article but do not replace its proofs.
\end{abstract}
\vfill
\small
\textbf{Review status.} This is an original research draft with human-readable
proofs and executable checks, not a peer-reviewed publication or a Lean/Rocq
formalization. The literature and repository comparison is targeted, not
exhaustive. No repository files or OEIS entries have been modified.

[Write note, batch 73O1: filed in ProveIt's research-report collection;
provenance and status in Section~\ref{bgg:sec:collection}.]

[Write note, batch 86A, 4 October 2026: the report now has two Parts.
Part~II (Sections~13--21, Appendices~C--D), from a later manuscript,
proves the negativity conjecture $r_k<0$ for every $k\ge3$ and the
large-order law of $r_k$; see Section~\ref{bgg:sg:sec:collection}.]
\end{titlepage}
\hypersetup{pageanchor=true}

\tableofcontents
\newpage

\part{Integrality, congruences, and inversion}
\label{bgg:part:arithmetic}

\section{The problem and the results}
\label{bgg:sec:introduction}

Define
\begin{equation}\label{bgg:eq:generating-functions}
 f_0(z)=\exp\!\left(\frac{z}{1-z}\right),\qquad
 f_1(z)=\exp\!\left(\frac1{1-z}\right)E_1\!\left(\frac1{1-z}\right),
\end{equation}
where $E_1(x)=\int_x^\infty e^{-u}u^{-1}\dif u$ for $x>0$, and write
\[
 a_n=[z^n]f_0(z),\qquad b_n=[z^n]f_1(z),\qquad \rho_n=a_nb_n.
\]
These sequences arise naturally from special functions and combinatorics.
For example, $n!a_n$ counts partitions of an $n$-element labelled set into
nonempty linearly ordered blocks, the classical sequence A000262
\cite{oeis262}. The second sequence is a recessive companion of the same
second-order recurrence. Their product cancels the opposite stretched
exponentials and exposes a more subtle arithmetic structure.

Brent, Glasser, and Guttmann proved that
\begin{equation}\label{bgg:eq:known-asymptotics}
 \rho_n\sim-\frac1{2n^{3/2}}\sum_{k\ge0}d_kn^{-k},\qquad d_0=1,
 \qquad d_k\in\Q.
\end{equation}
Here and throughout, an all-orders expansion is meant in the Poincar\'e sense:
for each fixed $M\ge0$, truncation after $d_M$ has normalized remainder
$O(n^{-M-1})$. We use the analytic existence statement from their Corollary~9;
we do not claim to rediscover it \cite{BGG}.

Their Conjecture~10 asserts that $r_k=64^kd_k$ is integral. Their Remark~11
records two additional observations for $3\le k\le1000$: negativity and the
congruence $r_k\equiv\binom{2k}{k}\pmod{32}$. Their proved Theorem~17 is the
weaker statement $k!r_k\in\Z$ \cite{BGG}. The OEIS entry A321941 still describes
integrality as conjectural in the entry inspected on 1 October 2026
\cite{oeis321941}. Its theorem numbers refer to an earlier preprint; our
numbers refer to the final 2019 journal article.

\begin{theorem}[Integrality and the full mod-$32$ statement]\label{bgg:thm:main}
For the coefficients in \eqref{bgg:eq:known-asymptotics},
\begin{equation}\label{bgg:eq:main-even}
 r_0=1,\qquad r_k=64^kd_k\in2\Z\quad(k\ge1).
\end{equation}
Moreover,
\begin{equation}\label{bgg:eq:main-congruence}
 r_k\equiv\binom{2k}{k}+16\ind_{\{1,2\}}(k)\pmod{32}
 \qquad(k\ge0).
\end{equation}
Thus the integrality conjecture and the congruence part of the stronger
observation are true at every index.
\end{theorem}

The exceptional corrections in \eqref{bgg:eq:main-congruence} are necessary:
$r_1=-14$ and $r_2=86$, whereas the corresponding central binomial
coefficients are $2$ and $6$.

\begin{theorem}[Denominators]\label{bgg:thm:denominators-summary}
The positive integers $B$ for which $B^kd_k\in\Z$ for every $k\ge0$ are
exactly the multiples of $64$. For every $k\ge1$,
\begin{equation}\label{bgg:eq:denominator-bound}
 \den(d_k)\mid 2^{6k-1}.
\end{equation}
If $s_2(k)$ denotes the number of ones in the binary expansion of $k$, then
\begin{equation}\label{bgg:eq:exact-denominator}
 \den(d_k)=2^{6k-s_2(k)}\qquad\text{whenever }1\le s_2(k)\le4.
\end{equation}
In particular, \eqref{bgg:eq:denominator-bound} is attained whenever $k$ is a
power of two.
\end{theorem}

Theorem~\ref{bgg:thm:main} is proved in Sections~\ref{bgg:sec:integrality}
and~\ref{bgg:sec:congruence}. Theorem~\ref{bgg:thm:denominators-summary} is proved in
Section~\ref{bgg:sec:denominators}. We also construct an integral polynomial family
in Section~\ref{bgg:sec:polynomials} and invert the asymptotic expansion in
Section~\ref{bgg:sec:inversion}.

\subsection{What is new, and what is not claimed}

The contribution is an arithmetic mechanism for this published conjecture,
not the invention of asymptotic series, Lagrange inversion, or products of
solutions of second-order equations. The crucial change is to use a
\emph{quadratic} product identity carrying the Wronskian, rather than only a
\emph{linear} recurrence for the product. In the new equation, the unknown
coefficient is multiplied by $2$, not by its index. Evenness then makes the
only division exact.

The proof also yields a polynomial deformation and congruences that are not
asserted in the cited source. These are presented as results established in
this draft; global literature priority has not been certified. A targeted
search for A321941 and for the published conjecture found the original paper
and OEIS entry, but no subsequent proof. A repository code search for
\texttt{A321941} returned no matches; this is not a proof of absence from every
file in the repository.

The separate assertion $r_k<0$ for all $k\ge3$ remains open in this work.
No convergence, optimal truncation theorem, or exponentially small completion
of \eqref{bgg:eq:known-asymptotics} is inferred from formal coefficient arithmetic.

\emph{[Write note, batch 86A, 4 October 2026.]} The assertion $r_k<0$ for
every $k\ge3$ is proved in Part~\ref{bgg:sg:part:sign}
(Theorem~\ref{bgg:sg:thm:sign}), by a different, analytic method; the
statements of this Part about it describe the work of 1 October 2026.
Convergence, optimal truncation and exponentially small completions
remain open.

\subsection{Place in the ProveIt collection}
\label{bgg:sec:collection}

\emph{[Write note, batch 73O1, 1 October 2026.]} This report is filed in
ProveIt's research-report collection, in the directory
\begin{center}\small
\path{SetTheory/Cardinals/docs/reports/congruences-and-valuations/}\\
\path{a321941-asymptotic-coefficient-integrality}.
\end{center}
Its README there is the guide to the files.

\emph{Provenance.} The report is built from one manuscript: manuscript~45
of batch~73 (cluster O1) of ProveIt's incoming-reports intake, the archive
\texttt{A321941\_Arithmetic\_Conjectures.zip}, which arrived in commit
\texttt{c79d64038} and was placed by commit \texttt{9df4ba51a}. Its pin is
the commit \repoSHA{} quoted in Section~\ref{bgg:sec:proveit}. The text is
printed as delivered. The write changed only the label names, which now
carry the prefix \texttt{bgg:}, and added the notes marked
\emph{[Write note]}. Nothing was merged into it; no other manuscript of
the batch treats this sequence.

\emph{Status.} The manuscript's author line is ``ChatGPT''. It is
unrefereed, and nothing in it is formalized in Lean, Rocq or any other
proof assistant. The intake reran the shipped verifier on a copy (exit
code~0; the coefficient and numerical tables byte-identical to the shipped
ones). Independently of the delivered code, it derived the
symmetric-square recurrence of $\rho_n$ from the recurrence of $a_n$,
solved it order by order for $r_0,\ldots,r_{10}$ (identical to the shipped
table), and confirmed the residues modulo~$32$ for $k\le10$ and
$\den(d_k)=2^{6k-s_2(k)}$ for $k\le8$. It did not re-derive every proof.
The statements about Brent--Glasser--Guttmann's paper (the range
$3\le k\le1000$ of their Remark~11, their theorem numbers) and about the
OEIS entry (its theorem numbering, and integrality still described as
conjectural) are as the manuscript inspected them on 1 October 2026; the
intake did not re-check them.

\emph{Shipped layout.} The verifier is \texttt{code/verify.py}; its
recorded outputs, the console log and the build receipt are in
\texttt{data/}. The paths printed in Section~\ref{bgg:sec:checks} and
Appendix~\ref{bgg:app:audit} are the delivered ones; the delivered README and PDF are not
shipped.

\emph{Formal neighbour.} The Lean module
\path{Analysis/FabiusFunction/Lean/FabiusFunction/QuadraticCoreCatalan.lean},
cited in Section~\ref{bgg:sec:proveit} and unchanged since the pin,
formalizes the Catalan closed form of $\binom{1/2}{n}$
(\texttt{Fabius.quadHalf\_rat}) and the Catalan convolution
(\texttt{Fabius.quadHalf\_antidiagonal}). These are an ingredient of the
fixed-point formulation \eqref{bgg:eq:catalan-square-root}, not of any
theorem of this report as stated; the power-series identity itself is not
formalized there, and this report gains no formal status from the module.

\emph{Neighbouring reports.} In the same collection,
\path{secant-number-periodicity} (OEIS A000364 modulo $m$) and
\path{iterated-series/compositional-tree-series-congruences} (integrality
and congruences of a formal fixed point) are arithmetic studies of other
sequences; there is no overlap. The negativity assertion $r_k<0$ for
$k\ge3$ remains open (question~1 of Section~\ref{bgg:sec:future}).

\emph{[Write note, batch 86A, 4 October 2026.]} The previous sentence is
superseded: Part~\ref{bgg:sg:part:sign} proves $r_k<0$ for every $k\ge3$,
with the large-order law of $r_k$; its provenance, status and notation
are in Section~\ref{bgg:sg:sec:collection}.

\section{The recurrence and its quadratic product invariant}
\label{bgg:sec:invariant}

\subsection{A unit-determinant normalization}

Put $G=eE_1(1)$. Direct differentiation of \eqref{bgg:eq:generating-functions}
gives
\begin{equation}\label{bgg:eq:odes}
 (1-z)^2f_0'(z)-f_0(z)=0,\qquad
 (1-z)^2f_1'(z)-f_1(z)=z-1.
\end{equation}
For the second identity, use $E_1'(x)=-e^{-x}/x$ and
$x=(1-z)^{-1}$. Coefficient extraction yields
\begin{equation}\label{bgg:eq:original-recurrence}
 (n+1)u_{n+1}-(2n+1)u_n+(n-1)u_{n-1}=0\qquad(n\ge2)
\end{equation}
for both $u=a$ and $u=b$. The initial values are
\[
 a_1=1,\quad a_2=\frac32,\qquad
 b_1=G-1,\quad b_2=\frac{3G-2}{2}.
\]
Set
\begin{equation}\label{bgg:eq:normalized-solutions}
 x_n=na_n,\qquad y_n=-nb_n,\qquad p_n=x_ny_n=-n^2\rho_n.
\end{equation}
Then both normalized solutions satisfy
\begin{equation}\label{bgg:eq:unimodular-recurrence}
 u_{n+1}=q_nu_n-u_{n-1},\qquad q_n=2+\frac1n\qquad(n\ge2).
\end{equation}
Their discrete Wronskian
\[
 W_n=x_ny_{n+1}-x_{n+1}y_n
\]
is constant: substituting \eqref{bgg:eq:unimodular-recurrence} shows
$W_n=W_{n-1}$. The initial values give
\[
 W_1=1(2-3G)-3(1-G)=-1.
\]
The exact value $W_n^2=1$ is the normalization needed below.

For completeness, $p_n>0$ for $n\ge1$. The coefficients $a_n$ are positive.
For the other factor, the integral formula
\begin{equation}\label{bgg:eq:b-integral}
 b_n=-\int_0^1 e^{-s/(1-s)}s^{n-1}\dif s\qquad(n\ge1)
\end{equation}
shows $b_n<0$. One derivation starts from
$e^xE_1(x)=\int_0^\infty e^{-u}(x+u)^{-1}\dif u$, extracts the coefficient
of $z^n$ at $x=(1-z)^{-1}$, and substitutes $s=u/(1+u)$.

\subsection{A general algebraic identity}

\begin{lemma}[Quadratic invariant]\label{bgg:lem:invariant}
Suppose $x_{n+1}=q_nx_n-x_{n-1}$ and
$y_{n+1}=q_ny_n-y_{n-1}$. Write $p_n=x_ny_n$ and
$W=x_ny_{n+1}-x_{n+1}y_n$. Then
\begin{equation}\label{bgg:eq:invariant}
 q_n^4p_n^2-2q_n^2p_n(p_{n+1}+p_{n-1})
 +(p_{n+1}-p_{n-1})^2=q_n^2W^2.
\end{equation}
This is a polynomial identity; no division by $p_n$ or $q_n$ is required.
\end{lemma}

\begin{proof}
Define $C_n=x_{n+1}y_n+x_ny_{n+1}$. The difference-of-squares identity gives
\[
 C_n^2-4p_np_{n+1}=W^2.
\]
Using $x_{n-1}=q_nx_n-x_{n+1}$ and the analogous identity for $y$ gives
\[
 p_{n-1}=q_n^2p_n-q_nC_n+p_{n+1},
\]
so $q_nC_n=q_n^2p_n+p_{n+1}-p_{n-1}$. Multiply the first equation by $q_n^2$
and substitute this expression. Expanding and collecting terms proves
\eqref{bgg:eq:invariant}.
\end{proof}

The Wronskian condition supplies the inhomogeneous constant on the right.
Discarding it would leave an unnecessarily ambiguous formal product equation.
In our application, \eqref{bgg:eq:invariant} holds for every integer $n\ge2$,
with $q_n=2+1/n$ and $W^2=1$.

\section{From the exact recurrence to a formal master equation}
\label{bgg:sec:master}

Let
\begin{equation}\label{bgg:eq:D-definition}
 D(h)=\sum_{k\ge0}d_kh^k\in\Q[[h]],\qquad D(0)=1.
\end{equation}
This is a formal series: we do not evaluate an infinite convergent sum at
$h=1/n$. Equation~\eqref{bgg:eq:known-asymptotics} says that
$p_n\sim\frac12\sqrt n\,D(1/n)$ in the all-orders sense.

For any formal series $D$, define the two shifted expressions
\begin{equation}\label{bgg:eq:shift-U}
 U_\pm(h)=(1\pm h)^{1/2}D\!\left(\frac h{1\pm h}\right).
\end{equation}
They encode $p_{n\pm1}$ after the common factor $\frac12\sqrt n$ is removed.
Substitution into \eqref{bgg:eq:invariant} gives, with $q=2+h$,
\begin{equation}\label{bgg:eq:unscaled-invariant}
 q^2D^2-2D(U_++U_-)+\frac{(U_+-U_-)^2}{q^2}=4h.
\end{equation}

\begin{lemma}[Justification of the formal identity]\label{bgg:lem:formal-bridge}
The coefficients from the published expansion satisfy
\eqref{bgg:eq:unscaled-invariant} as an identity in $\Q[[h]]$.
\end{lemma}

\begin{proof}
Fix an arbitrary finite target order. Expand $p_n$ and $p_{n\pm1}$ to
sufficiently many additional orders in the exact identity
\eqref{bgg:eq:invariant}. Shifts by a fixed integer preserve Poincar\'e
expansions; the substitution $1/(n\pm1)=h/(1\pm h)$ and the binomial expansion
of $(1\pm h)^{1/2}$ give precisely \eqref{bgg:eq:shift-U}. Products and division
by the unit $q^2=(2+h)^2$ preserve finite-order expansions. After removing the
common factor $n/4$, the exact identity gives
\eqref{bgg:eq:unscaled-invariant} through the chosen order. A first nonzero formal
coefficient would contradict this conclusion at its order. Since the target
order was arbitrary, every coefficient agrees.
\end{proof}

The even and odd parts of the shift make cancellation explicit. Put
\begin{equation}\label{bgg:eq:unscaled-operators}
 E(D)=\frac{U_++U_--2D}{h^2},\qquad
 O(D)=\frac{U_+-U_-}{h}.
\end{equation}
Both quotients are formal series, because the indicated numerators have the
required orders. Since
\[
 U_\pm=\sum_{j\ge0}d_jh^j(1\pm h)^{1/2-j},
\]
their coefficients can be written explicitly using half-integral binomial
coefficients. Inserting \eqref{bgg:eq:unscaled-operators} into
\eqref{bgg:eq:unscaled-invariant}, and using $q^2-4=4h+h^2$, gives
\begin{equation}\label{bgg:eq:D-master}
 D^2=1-\frac h4D^2+\frac h2DE(D)
       -\frac{h}{4(2+h)^2}O(D)^2.
\end{equation}
The cancellation of one factor $h$ is legitimate in $\Q[[h]]$.

Now scale
\begin{equation}\label{bgg:eq:R-definition}
 R(z)=D(64z)=\sum_{k\ge0}r_kz^k.
\end{equation}
For $j,\ell\ge0$, define rational weights
\begin{equation}\label{bgg:eq:weights}
 e_{j,\ell}=8\cdot64^{2\ell}\binom{1/2-j}{2\ell+2},\qquad
 o_{j,\ell}=2\cdot64^{2\ell}\binom{1/2-j}{2\ell+1}.
\end{equation}
For a series $F(z)=\sum_{j\ge0}f_jz^j$, set
\begin{equation}\label{bgg:eq:scaled-operators}
 \Eop F=\sum_{j,\ell\ge0}e_{j,\ell}f_jz^{j+2\ell},\qquad
 \Oop F=\sum_{j,\ell\ge0}o_{j,\ell}f_jz^{j+2\ell}.
\end{equation}
Each coefficient is a finite sum. Directly from the definitions,
$\Eop R=4E(D)(64z)$ and $\Oop R=O(D)(64z)$.

\begin{proposition}[The scaled master equation]\label{bgg:prop:master}
The series $R$ satisfies
\begin{equation}\label{bgg:eq:master}
 \boxed{\displaystyle
 R^2=1+4z\left(-4R^2+2R\Eop R
              -\frac{(\Oop R)^2}{(1+32z)^2}\right).}
\end{equation}
Conversely, among formal series in $\Q[[z]]$ with constant term $1$,
\eqref{bgg:eq:master} has at most one solution.
\end{proposition}

\begin{proof}
The identity is the substitution $h=64z$ in \eqref{bgg:eq:D-master}. For
uniqueness, the coefficient of $z^k$ on the left is $2r_k$ plus an expression
in earlier coefficients. The coefficient of $z^k$ on the right depends only
on coefficients through $r_{k-1}$: there is an external factor $z$, and the
operators never decrease degree. Therefore $r_k$ is determined recursively.
\end{proof}

\section{Integral shift operators and the integrality proof}
\label{bgg:sec:integrality}

\subsection{The exact powers of two in the shift weights}

\begin{lemma}[Half-binomial arithmetic]\label{bgg:lem:half-binomial}
For integers $j\ge0$ and $m\ge1$, the number
$\binom{1/2-j}{m}$ is dyadic, and
\begin{equation}\label{bgg:eq:half-valuation}
 v_2\!\binom{1/2-j}{m}=-m-v_2(m!)=-2m+s_2(m).
\end{equation}
Consequently, all weights in \eqref{bgg:eq:weights} are integers. More precisely,
\begin{align}
 e_{j,0}&=4j^2-1,&o_{j,0}&=1-2j,\label{bgg:eq:diagonal-weights}\\
 256&\mid e_{j,\ell},&512&\mid o_{j,\ell}\qquad(\ell\ge1).
 \label{bgg:eq:offdiagonal-divisibility}
\end{align}
\end{lemma}

\begin{proof}
The elementary formula
\[
 \binom{1/2}{m}=\frac{(-1)^{m-1}\Cat_{m-1}}{2^{2m-1}}\qquad(m\ge1),
 \qquad \Cat_r=\frac1{r+1}\binom{2r}{r}\in\Z,
\]
shows that the coefficients of $(1+x)^{1/2}$ are dyadic. Multiplication by
$(1+x)^{-j}\in\Z[[x]]$ proves the same for
$\binom{1/2-j}{m}$. Alternatively, this is the finite Vandermonde convolution.
In the product
\[
 \binom{1/2-j}{m}
 =\frac{\prod_{i=0}^{m-1}(1-2j-2i)}{2^m m!},
\]
the numerator is odd. This proves the first equality in
\eqref{bgg:eq:half-valuation}. The identity
$v_2(m!)=m-s_2(m)$ gives the second.

For the even and odd weights, respectively, substitution yields
\begin{equation}\label{bgg:eq:weight-valuations}
 v_2(e_{j,\ell})=8\ell-1+s_2(\ell+1),\qquad
 v_2(o_{j,\ell})=8\ell+s_2(\ell).
\end{equation}
For $\ell=0$, both valuations are zero, and direct calculation gives
\eqref{bgg:eq:diagonal-weights}. For $\ell\ge1$, the valuations are at least $8$
and $9$, respectively. Dyadicity together with these nonnegative valuations
proves integrality, not merely the absence of a factor $2$ in a denominator.
\end{proof}

Writing $\theta=z\frac{\mathrm d}{\mathrm dz}$, the lemma implies, for every
$F\in\Z[[z]]$,
\begin{equation}\label{bgg:eq:operator-congruences}
 \Eop F-(4\theta^2-1)F\in256z^2\Z[[z]],\qquad
 \Oop F-(1-2\theta)F\in512z^2\Z[[z]].
\end{equation}
These stronger divisibilities will be used for the congruence theorem.

\subsection{The triangular recurrence}

Define
\begin{equation}\label{bgg:eq:H-definition}
 \Hop(F)=-4F^2+2F\Eop F-\left(\frac{\Oop F}{1+32z}\right)^2.
\end{equation}
Lemma~\ref{bgg:lem:half-binomial} shows that $\Hop$ maps integral series to
integral series. Let $R_{<k}=\sum_{j=0}^{k-1}r_jz^j$. Extracting $z^k$ from
\eqref{bgg:eq:master} gives
\begin{equation}\label{bgg:eq:triangular-recurrence}
 \boxed{\displaystyle
 r_k=2[z^{k-1}]\Hop(R_{<k})
     -\frac12\sum_{j=1}^{k-1}r_jr_{k-j}.}
\end{equation}
This expression depends only on already determined coefficients.

\begin{proof}[Proof of the integrality assertion in Theorem~\ref{bgg:thm:main}]
Start with $r_0=1$. Suppose $r_j\in2\Z$ for $1\le j<k$. Then
$[z^{k-1}]\Hop(R_{<k})$ is an integer. Every summand in the convolution in
\eqref{bgg:eq:triangular-recurrence} is divisible by $4$, so half of the whole sum
is an even integer. Both terms on the right are therefore even integers.
This proves $r_k\in2\Z$ by induction. Proposition~\ref{bgg:prop:master} and
Lemma~\ref{bgg:lem:formal-bridge} identify this recursively constructed series
with the coefficients of the analytic expansion.
\end{proof}

The original index division has disappeared. This is why the quadratic
identity proves more than the published linear recurrence, whose direct
arithmetic consequence was $k!r_k\in\Z$. The proof establishes actual
integrality of $r_k$; it does not assume the finite OEIS values are integers
and extrapolate.

The first steps illustrate the mechanism. At constant order,
$\Eop1=-1+O(z^2)$ and $\Oop1=1+O(z^2)$, so
$\Hop(1)=-7+O(z)$. Thus $r_1=2(-7)=-14$.
Subsequent applications give $r_2=86$, $r_3=-3660$, and the table in
Section~\ref{bgg:sec:checks}.

\subsection{A Catalan fixed-point interpretation}

There is also a useful division-free formulation. The formal identity
\begin{equation}\label{bgg:eq:catalan-square-root}
 \sqrt{1+4X}=1+2\sum_{m\ge1}(-1)^{m-1}\Cat_{m-1}X^m
\end{equation}
has integer coefficients. Thus the map
\[
 \Phi(F)=1+2\sum_{m\ge1}(-1)^{m-1}\Cat_{m-1}
                 \bigl(z\Hop(F)\bigr)^m
\]
maps $\Z[[z]]$ into $1+2z\Z[[z]]$. Because the operators in
$\Hop$ do not decrease degree, agreement of $F$ and $G$ through degree $k-1$
implies agreement of $\Phi(F)$ and $\Phi(G)$ through degree $k$.
Successive substitution therefore determines a unique formal fixed point.
Over $\Q[[z]]$, squaring identifies its equation with \eqref{bgg:eq:master}.

This interpretation explains why Catalan arithmetic is relevant. It also
suggests a particularly clean formalization route: construct the fixed point
by finite coefficient truncations and use the Catalan convolution, rather
than formal square-root analysis.

\section{An integral polynomial deformation}
\label{bgg:sec:polynomials}

The same proof works with an indeterminate parameter. This is more than an
integer sequence coincidence: the coefficient arithmetic survives a
one-parameter deformation of the discrete equation.

\subsection{Definition and integrality}

Let $t$ be an indeterminate. For $F=\sum f_j(t)z^j$, define
\begin{equation}\label{bgg:eq:parameter-operators}
 \Eop_tF=\sum_{j,\ell\ge0}e_{j,\ell}t^\ell f_j(t)z^{j+2\ell},\qquad
 \Oop_tF=\sum_{j,\ell\ge0}o_{j,\ell}t^\ell f_j(t)z^{j+2\ell}.
\end{equation}
These are well-defined, degree-nondecreasing linear operators on
$\Z[t][[z]]$. Put
\begin{equation}\label{bgg:eq:H-parameter}
 \Hop_t(F)=-4tF^2+2F\Eop_tF
          -\left(\frac{\Oop_tF}{1+32tz}\right)^2.
\end{equation}

\begin{theorem}[Polynomial master series]\label{bgg:thm:polynomial-integrality}
There is a unique series $R_t\in1+z\Q[t][[z]]$ satisfying
\begin{equation}\label{bgg:eq:parameter-master}
 R_t^2=1+4z\Hop_t(R_t).
\end{equation}
Writing $R_t(z)=\sum_{k\ge0}P_k(t)z^k$, one has
\begin{equation}\label{bgg:eq:polynomial-integrality}
 P_0(t)=1,\qquad P_k(t)\in2\Z[t]\quad(k\ge1),\qquad P_k(1)=r_k.
\end{equation}
Every $P_k$ has degree exactly $k$, with
\begin{equation}\label{bgg:eq:leading-polynomial}
 [t^k]P_k(t)=(-4)^k\binom{2k}{k}.
\end{equation}
\end{theorem}

\begin{proof}
The recurrence \eqref{bgg:eq:triangular-recurrence}, with $\Hop$ replaced by
$\Hop_t$, determines each coefficient from its predecessors. The evenness
induction works coefficientwise over $\Z[t]$ and proves existence and
\eqref{bgg:eq:polynomial-integrality}. Uniqueness over $\Q[t]$ follows from the
coefficient $2P_k$ on the left. At $t=1$, the equation is
\eqref{bgg:eq:master}, proving $P_k(1)=r_k$.

Inductively, assume $\deg P_j\le j$. In the coefficient of $z^m$ in
$\Eop_tR_t$ or $\Oop_tR_t$, a summand with $j+2\ell=m$ has degree at most
$j+\ell\le m$. The same degree bound is preserved by multiplication and by
$(1+32tz)^{-1}$. Consequently, the recurrence gives $\deg P_k\le k$.

In total bidegree $t^kz^k$, the terms $8zR_t\Eop_tR_t$ and
$-4z(\Oop_tR_t/(1+32tz))^2$ have insufficient $t$-degree. If
$S(w)=\sum_{k\ge0}[t^k]P_k(t)w^k$, the diagonal terms of
\eqref{bgg:eq:parameter-master} therefore satisfy
\[
 S(w)^2=1-16wS(w)^2.
\]
The branch with constant term $1$ is $S(w)=(1+16w)^{-1/2}$, whose $k$th
coefficient is $(-4)^k\binom{2k}{k}\ne0$. This proves both the degree and
\eqref{bgg:eq:leading-polynomial}.
\end{proof}

The first polynomials are
\begin{align*}
 P_1(t)&=-8t-6,\\
 P_2(t)&=96t^2+80t-90,\\
 P_3(t)&=-1280t^3+2240t^2+1680t-6300,\\
 P_4(t)&=17920t^4-161280t^3-57792t^2+151200t-992250.
\end{align*}
The alternating and mixed signs already show why integrality alone does not
prove the separate negativity assertion at $t=1$.

\subsection{The constant edge has a closed form}

\begin{theorem}[Constant coefficients]\label{bgg:thm:constant-edge}
For every $k\ge1$,
\begin{equation}\label{bgg:eq:constant-polynomial}
 P_k(0)=-\binom{2k}{k}(2k-3)!!(2k+1)!!,
\end{equation}
where $(-1)!!=1$. In particular, this edge of the polynomial family is
strictly negative at every positive index.
\end{theorem}

\begin{proof}
Set $S(z)=R_0(z)$. At $t=0$, only the diagonal operator weights survive:
$\Eop_0S=(4\theta^2-1)S$ and $\Oop_0S=(1-2\theta)S$. Hence
\begin{equation}\label{bgg:eq:constant-edge-nonlinear}
 S^2=1+4z\bigl(-3S^2+4S\theta S+8S\theta^2S-4(\theta S)^2\bigr).
\end{equation}
Apply $\theta$ to both sides. If the parenthesized expression is denoted by
$F$, its derivative satisfies
\[
 F+\theta F
 =S\bigl(-3S-2\theta S+12\theta^2S+8\theta^3S\bigr).
\]
Cancel the unit $S$ to obtain
\begin{equation}\label{bgg:eq:constant-edge-linear}
 \theta S=2z(2\theta+3)(2\theta-1)(2\theta+1)S.
\end{equation}
Coefficient extraction gives
\begin{equation}\label{bgg:eq:constant-edge-recurrence}
 kP_k(0)=2(2k+1)(2k-3)(2k-1)P_{k-1}(0).
\end{equation}
The expression on the right of \eqref{bgg:eq:constant-polynomial} has initial
value $-6$ at $k=1$ and the same consecutive ratio. Since $P_0(0)=1$,
\eqref{bgg:eq:constant-edge-recurrence} proves the formula by induction.
\end{proof}

Thus the extreme coefficients of each polynomial are explicit, although the
interior coefficients need not have one sign. Formula
\eqref{bgg:eq:constant-edge-linear} is a useful linear equation at this special
parameter, not a replacement for the nonlinear integrality argument at
$t=1$.

\emph{[Write note, batch 86A, 4 October 2026.]} The negative edge
$P_k(0)=-S_k$ is the leading asymptotic scale of $P_k(1)=r_k$:
Theorem~\ref{bgg:sg:thm:largeorder} of Part~\ref{bgg:sg:part:sign} gives
$r_k/P_k(0)=1-1/(3k)+O(k^{-2})$, with an expansion to every order.

\subsection{Why this deformation is natural}

Consider the recurrence $u_{n+1}=(2+\lambda/n)u_n-u_{n-1}$ over a field
containing a nonzero parameter $\lambda$ and a square root of $\lambda$.
Seek a formal solution of the quadratic identity, with Wronskian square $1$,
of the form
\[
 p_n\sim\frac{\sqrt n}{2\sqrt\lambda}D_\lambda(1/n),
 \qquad D_\lambda(0)=1.
\]
Here ``formal product'' means the normalized solution of the quadratic
identity; analytic existence for all parameters is not assumed. Repeating
the calculation of Section~\ref{bgg:sec:master} gives
\begin{equation}\label{bgg:eq:lambda-master}
 D_\lambda^2=1-\frac{\lambda h}{4}D_\lambda^2
 +\frac{h}{2\lambda}D_\lambda E(D_\lambda)
 -\frac{h}{4\lambda(2+\lambda h)^2}O(D_\lambda)^2.
\end{equation}
Set $h=64\lambda z$ and $t=\lambda^2$. The even and odd shift weights then
become exactly \eqref{bgg:eq:parameter-operators}, and
\eqref{bgg:eq:lambda-master} becomes \eqref{bgg:eq:parameter-master}. Therefore
\begin{equation}\label{bgg:eq:lambda-coefficients}
 [h^k]D_\lambda(h)=\frac{P_k(\lambda^2)}{(64\lambda)^k}.
\end{equation}
This is an unconditional statement about the unique normalized formal
solution. Whenever an actual normalized product has an all-orders expansion,
the formal-bridge argument identifies its coefficients with
\eqref{bgg:eq:lambda-coefficients}. The case $\lambda=1$ is established by the
published analytic theorem. The specialization $t=0$ in
Theorem~\ref{bgg:thm:constant-edge} is a formal polynomial specialization; it is
not a claimed uniform analytic limit as $\lambda\to0$.

\section{The congruence modulo 32}
\label{bgg:sec:congruence}

The large off-diagonal powers of two in Lemma~\ref{bgg:lem:half-binomial} allow a
short reduction of the entire master equation. We prove the stronger
polynomial statement first.

\begin{theorem}[Polynomial congruence]\label{bgg:thm:polynomial-congruence}
For every $k\ge0$, coefficientwise in $\Z[t]$,
\begin{equation}\label{bgg:eq:polynomial-congruence}
 P_k(t)\equiv(-1)^k(4t+3)^k\binom{2k}{k}
             +16\ind_{\{2\}}(k)\pmod{32}.
\end{equation}
Equivalently,
\begin{equation}\label{bgg:eq:series-congruence}
 R_t(z)\equiv(1+(16t+12)z)^{-1/2}+16z^2\pmod{32}.
\end{equation}
\end{theorem}

\begin{proof}
Write $R=R_t$ temporarily. The operator congruences
\eqref{bgg:eq:operator-congruences} remain valid over $\Z[t]$, and
$(1+32tz)^{-2}\equiv1\pmod{64}$. Reducing
\eqref{bgg:eq:parameter-master} modulo $64$ gives
\begin{equation}\label{bgg:eq:mod64-intermediate}
 R^2\equiv1+4z\bigl(-(4t+3)R^2+4R\theta R
                 +8R\theta^2R-4(\theta R)^2\bigr)\pmod{64}.
\end{equation}
Theorem~\ref{bgg:thm:polynomial-integrality} shows that both $\theta R$ and
$\theta^2R$ are even series. Hence the last two terms, after multiplication
by $4z$, vanish modulo $64$. With $\alpha=16t+12$, we obtain
\begin{equation}\label{bgg:eq:reduced-equation}
 (1+\alpha z)R^2\equiv1+16zR\theta R\pmod{64}.
\end{equation}

Let $C=(1+\alpha z)^{-1/2}$. Its coefficients are
$(-1)^k(4t+3)^k\binom{2k}{k}$, so $C\in1+2z\Z[t][[z]]$. Define
$S=C+16z^2$. Since $C\equiv1\pmod2$ and $\alpha\equiv0\pmod4$,
\begin{align*}
 (1+\alpha z)S^2&\equiv1+32z^2\pmod{64},\\
 16zS\theta S&\equiv16zC\theta C
       =-\frac{8\alpha z^2}{(1+\alpha z)^2}
       \equiv32z^2\pmod{64}.
\end{align*}
For the last congruence, $-8\alpha\equiv32\pmod{64}$, and
$(1+\alpha z)^{-2}\equiv1\pmod2$. Thus $S$ satisfies the same reduced
equation as $R$.

We now justify the loss of exactly one power of two when taking the square
root; no division by $2$ in $\Z/64\Z$ is being assumed. Both $R$ and $S$
have constant term $1$ and even higher coefficients. Suppose their
coefficients below degree $k$ agree modulo $32$. In degree $k$,
$R^2-S^2$ is $2([z^k]R-[z^k]S)$ modulo $64$: every remaining product
contains an even nonconstant coefficient and a difference divisible by $32$.
The term $\alpha z(R^2-S^2)$ and the difference
$16z(R\theta R-S\theta S)$ depend only on lower degrees and vanish modulo
$64$. Subtracting their two instances of \eqref{bgg:eq:reduced-equation}
therefore gives
\[
 2([z^k]R-[z^k]S)\equiv0\pmod{64}.
\]
Thus $[z^k]R\equiv[z^k]S\pmod{32}$. Induction proves
\eqref{bgg:eq:series-congruence}, and coefficient extraction gives
\eqref{bgg:eq:polynomial-congruence}.
\end{proof}

\begin{proof}[Proof of the congruence assertion in Theorem~\ref{bgg:thm:main}]
Specialize \eqref{bgg:eq:series-congruence} to $t=1$:
\[
 R(z)\equiv(1+28z)^{-1/2}+16z^2\pmod{32}.
\]
Factor $1+28z=(1-4z)(1+32z/(1-4z))$. The binomial expansion gives
\[
 (1+28z)^{-1/2}
 \equiv(1-4z)^{-1/2}-16z(1-4z)^{-3/2}\pmod{32}.
\]
All terms of second and higher degree in $32z/(1-4z)$ vanish modulo $32$;
this follows directly from their half-binomial valuations. The nonconstant
coefficients of $(1-4z)^{-3/2}$ are even: its $k$th coefficient is
$(2k+1)\binom{2k}{k}$. Consequently,
\begin{equation}\label{bgg:eq:final-R-congruence}
 R(z)\equiv(1-4z)^{-1/2}+16z+16z^2\pmod{32}.
\end{equation}
Using $[z^k](1-4z)^{-1/2}=\binom{2k}{k}$ proves
\eqref{bgg:eq:main-congruence} at all indices, including the exceptions.
\end{proof}

The proof demonstrates why a central binomial sequence appears: modulo a
small power of two, the nonlocal shift equation collapses to a differential
quadratic equation whose solution is an inverse square root, up to two
low-degree corrections.

\section{Minimal scaling and exact denominators}
\label{bgg:sec:denominators}

\begin{proof}[Proof of Theorem~\ref{bgg:thm:denominators-summary}]
Since $r_k$ is even for $k\ge1$, the denominator of $d_k=r_k/2^{6k}$ divides
$2^{6k-1}$. The computed value
\[
 d_2=\frac{43}{2048}=\frac{43}{2^{11}}
\]
shows that $B^2d_2\in\Z$ forces $2v_2(B)\ge11$, hence $v_2(B)\ge6$.
Thus every permissible positive $B$ is divisible by $64$. Conversely,
$B=64m$ gives $B^kd_k=m^kr_k\in\Z$, proving the exact classification of
scales.

The factorial identity $v_2(k!)=k-s_2(k)$ gives
\begin{equation}\label{bgg:eq:central-valuation}
 v_2\binom{2k}{k}
 =(2k-s_2(2k))-2(k-s_2(k))=s_2(k).
\end{equation}
For $k\ge3$ with $s_2(k)\le4$, the congruence
$r_k\equiv\binom{2k}{k}\pmod{32}$ forces
$v_2(r_k)=s_2(k)$. Indeed, adding a multiple of $2^5$ cannot change a
valuation strictly less than $5$. The indices $k=1,2$ are checked directly:
$r_1=-14$ and $r_2=86$ both have valuation $1=s_2(k)$. Cancelling the exact
power of two in $d_k=r_k/2^{6k}$ proves
\eqref{bgg:eq:exact-denominator}. For $k=2^m$, $s_2(k)=1$, so the bound is
attained at every such index.
\end{proof}

This separates two notions of optimality. The base $64$ is the smallest
integer exponential scale making every coefficient integral. The additional
factor $2$ available at every positive index improves the denominator bound,
but cannot reduce that exponential base. The equality at infinitely many
powers of two shows that the improved bound is not merely an artifact of the
proof.

For indices with at least five binary ones, the mod-$32$ congruence only
forces $32\mid r_k$ and does not determine the exact valuation. In particular,
we do not extrapolate \eqref{bgg:eq:exact-denominator} to all indices.

\section{All-orders inversion of the product expansion}
\label{bgg:sec:inversion}

The arithmetic series also supports an explicit inverse expansion. It is
important to specify what is being inverted: a continuous global inverse of
an interpolation of the sequence has not been chosen.

Define the positive real sequence
\begin{equation}\label{bgg:eq:y-definition}
 y_n=(-2\rho_n)^{-2/3}.
\end{equation}
Positivity follows from \eqref{bgg:eq:b-integral}. The known product expansion
gives
\begin{equation}\label{bgg:eq:forward-inversion}
 y_n\sim nD(1/n)^{-2/3},\qquad y_n\sim n.
\end{equation}

\begin{theorem}[Inverse expansion]\label{bgg:thm:inverse}
Define
\begin{equation}\label{bgg:eq:inverse-coefficients}
 \beta_0=\frac23d_1,\qquad
 \beta_k=-\frac1k[h^{k+1}]D(h)^{-2k/3}\quad(k\ge1).
\end{equation}
Then, for each fixed $M\ge0$,
\begin{equation}\label{bgg:eq:inverse-expansion}
 n=y_n+\beta_0+\sum_{k=1}^{M}\beta_ky_n^{-k}
      +O(y_n^{-M-1})\qquad(n\to\infty).
\end{equation}
In particular,
\begin{align}\label{bgg:eq:displayed-inverse}
 n={}&y_n-\frac7{48}-\frac{29}{2304y_n}
       -\frac{3719}{331776y_n^2}-\frac{1565}{32768y_n^3}\notag\\
    &-\frac{36355159}{191102976y_n^4}+O(y_n^{-5}).
\end{align}
Moreover,
\begin{equation}\label{bgg:eq:inverse-arithmetic}
 64^{k+1}\beta_k\in\Z[1/3]\qquad(k\ge0).
\end{equation}
Thus every inverse coefficient has denominator supported only at the primes
$2$ and $3$, and after the displayed normalization only $3$ remains.
\end{theorem}

\begin{proof}
Let $h=1/n$ and formally put $u=1/y$. The forward relation becomes
\[
 u=hD(h)^{2/3},\qquad h=u\psi(h),\qquad\psi(h)=D(h)^{-2/3}.
\]
The series $u=h+O(h^2)$ has a unique formal compositional inverse. For $k\ge1$,
Lagrange inversion applied to the Laurent function $h^{-1}$ gives
\[
 [u^k]\frac1{h(u)}=-\frac1k[h^{k+1}]\psi(h)^k,
\]
which is \eqref{bgg:eq:inverse-coefficients}. This Laurent version follows, for
example, by changing variables in the coefficient residue and integrating a
formal derivative; it has the same finite coefficient meaning as the usual
power-series version. The constant term is obtained separately from
$h=u+\psi'(0)u^2+O(u^3)$: it is $-\psi'(0)=2d_1/3$.

For the asymptotic assertion, use the published expansion to enough fixed
orders and form \eqref{bgg:eq:forward-inversion} by finite-order Taylor
expansion around $1$. Substitute the truncated forward expansion into the
formal inverse. All coefficients through $y^{-M}$ cancel by construction.
Since $y_n\sim n$, the uncancelled terms and transported remainders are
$O(y_n^{-M-1})$. No global invertibility or monotonicity assumption is needed
for this composition along the sequence. Extracting the first coefficients
gives \eqref{bgg:eq:displayed-inverse}.

For arithmetic, introduce $X=64n$, $Y=64y$ and $z=1/X$. The normalized
forward relation is
\[
 Y=XR(1/X)^{-2/3},\qquad Y^{-1}=zR(z)^{2/3}.
\]
Because $R\in\Z[[z]]$, the unit $F=R^{2/3}$ belongs to
$\Z[1/3][[z]]$: solve $F^3=R^2$ coefficient by coefficient, dividing only
by $3$. A monic series $zF(z)=z+O(z^2)$ has a compositional inverse over the
same ring, and reciprocating that inverse gives
\[
 X(Y)\in Y+\Z[1/3][[Y^{-1}]].
\]
Since the coefficient of $Y^{-k}$ in $X(Y)-Y$ is
$64^{k+1}\beta_k$, this proves \eqref{bgg:eq:inverse-arithmetic}.
\end{proof}

\subsection{The boundary between an expansion and a transseries}

Theorem~\ref{bgg:thm:inverse} gives the complete algebraic sector of an inverse
asymptotic description. It does not specify exponentially small corrections.
Two functions differing by a term smaller than every inverse power have the
same formal series, so coefficient arithmetic cannot recover such a term.
This distinction is especially relevant here: the individual factors have
opposite stretched-exponential behavior, while their product has an
ordinary inverse-power expansion.

Similarly, an $O(y_n^{-M-1})$ statement supplies no numerical error constant
at a specified $n$. Rigorous rounding of an inverse estimate to an integer
would require an effective remainder bound. The numerical checks below
illustrate the expansion but are not interval error certificates.

\section{Reproducible exact and numerical checks}
\label{bgg:sec:checks}

\emph{[Write note, batch 73O1.]} In the collection the program is
\texttt{code/verify.py} and the files it writes are shipped in
\texttt{data/}. Its default output directory is its own, \texttt{code/};
pass \texttt{--out} to write elsewhere
(Appendix~\ref{bgg:app:audit}).

The accompanying program \texttt{verify.py} uses two different coefficient
recurrences. The new method computes $r_k$ using integer arithmetic and the
quadratic master equation. The independent method uses rational arithmetic
and the published linear coefficient recurrence. Agreement is a useful check
on normalization, signs, and indexing; the all-orders proofs above do not
depend on how far the calculations are run.

\subsection{An explicit implementation form}

Let
\[
 R=\sum r_nz^n,\qquad \Eop_tR=\sum E_nz^n,\qquad
 \Oop_tR=\sum O_nz^n,\qquad
 \frac{\Oop_tR}{1+32tz}=\sum T_nz^n.
\]
Starting with $r_0=1$, $E_0=-1$, $O_0=T_0=1$, the integer recurrence is
\begin{align}
 H_{n-1}&=-4t\sum_{i=0}^{n-1}r_ir_{n-1-i}
          +2\sum_{i=0}^{n-1}r_iE_{n-1-i}
          -\sum_{i=0}^{n-1}T_iT_{n-1-i},\label{bgg:eq:algorithm-H}\\
 r_n&=2H_{n-1}-\frac12\sum_{i=1}^{n-1}r_ir_{n-i},\label{bgg:eq:algorithm-r}\\
 E_n&=\sum_{\ell=0}^{\lfloor n/2\rfloor}
        e_{n-2\ell,\ell}t^\ell r_{n-2\ell},\qquad
 O_n=\sum_{\ell=0}^{\lfloor n/2\rfloor}
        o_{n-2\ell,\ell}t^\ell r_{n-2\ell},\label{bgg:eq:algorithm-EO}\\
 T_n&=O_n-32tT_{n-1}.\label{bgg:eq:algorithm-T}
\end{align}
The order matters: compute $r_n$ from earlier arrays, then append
$E_n,O_n,T_n$. The division by two in \eqref{bgg:eq:algorithm-r} is always exact,
and its quotient is even. The code asserts this at each step.

For comparison, the independent recurrence from \cite[Lemma~14]{BGG} is
\begin{equation}\label{bgg:eq:independent-recurrence}
 8k d_k=-\sum_{j=0}^{k-1}d_j[h^{k+2-j}]
 \left(\frac{B(h)}{(1-h)^{j+1/2}}+
       \frac{C(h)}{(1-2h)^{j+1/2}}+
       \frac{J(h)}{(1-3h)^{j+1/2}}\right),
\end{equation}
where
\[
 \begin{aligned}
 B(h)&=-6+13h-7h^2+h^3,\\
 C(h)&=6-19h+17h^2-3h^3,\\
 J(h)&=-2+11h-17h^2+6h^3.
 \end{aligned}
\]
The source calls the last polynomial $D$; we rename it $J$ to avoid collision
with the asymptotic series. This second implementation does not use
\eqref{bgg:eq:algorithm-r} or its integrality conclusion.

\subsection{Initial coefficients}

\begin{table}[ht]
\centering
\caption{Exact initial coefficients. The normalized product is
$-2n^{3/2}\rho_n\sim\sum d_kn^{-k}$.}
\label{bgg:tab:coefficients}
\begin{tabular}{r r r}
\toprule
$k$ & $r_k$ & $d_k$ in lowest terms\\
\midrule
0 & $1$ & $1$\\
1 & $-14$ & $-7/32$\\
2 & $86$ & $43/2048$\\
3 & $-3660$ & $-915/65536$\\
4 & $-1042202$ & $-521101/8388608$\\
5 & $-247948260$ & $-61987065/268435456$\\
6 & $-108448540420$ & $-27112135105/17179869184$\\
\bottomrule
\end{tabular}
\end{table}

The archived run used Python 3.13.5, SymPy 1.14.0, and mpmath 1.3.0. It
verified integrality/evenness, the mod-$32$ formula, and all applicable
valuation formulas through $k=256$; independent agreement with
\eqref{bgg:eq:independent-recurrence} through $k=128$; all thirteen displayed
OEIS terms; the polynomial identities through degree $8$; and the parameter
congruence through $k=48$ at $t=-3,-1,0,2,5$. It also checked the bilinear
invariant symbolically and composed the displayed inverse series directly.

The check that $r_k<0$ for $3\le k\le256$ is explicitly logged as an
\emph{observation}, not a proof. It does not extend the range already reported
in the 2019 paper and is not presented as a computational record.
\emph{[Write note, batch 86A: the infinite assertion is proved in
Part~\ref{bgg:sg:part:sign}, Theorem~\ref{bgg:sg:thm:sign}; this finite
check remains an observation of this Part.]}

\subsection{Numerical asymptotic checks}

For numerical evaluation we used the exact special-function expression
\begin{equation}\label{bgg:eq:kummer-numerical}
 \rho_n=-e^{-1}\Gamma(n)M(n+1,2,1)U(n,0,1),
\end{equation}
with the standard Kummer functions, as in \cite{BGG}. This avoids computing
the recessive $b_n$ by unstable forward propagation from nearly cancelling
initial data. Table~\ref{bgg:tab:numerical} shows the signed errors from the
published product and our inverse expansion.

\begin{table}[ht]
\centering
\caption{Numerical diagnostics, not rigorous remainder bounds. Product error:
$-2n^{3/2}\rho_n-\sum_{k=0}^{5}d_kn^{-k}$. Inverse error: the right side of
\eqref{bgg:eq:displayed-inverse}, without its $O$ term, minus $n$.}
\label{bgg:tab:numerical}
\begin{tabular}{r r r}
\toprule
$n$ & Product error & Inverse error\\
\midrule
25 & $-1.16668669188\times10^{-8}$ & $2.14328374121\times10^{-7}$\\
100 & $-1.75696329807\times10^{-12}$ & $1.34214874476\times10^{-10}$\\
400 & $-3.95020779621\times10^{-16}$ & $1.21710421134\times10^{-13}$\\
\bottomrule
\end{tabular}
\end{table}

The recorded numerical strings were unchanged when working precision was
increased from $80$ to $120$ decimal digits. This tests numerical stability
of the reported digits but is not interval arithmetic. Exact assertions are
separate from these numerical diagnostics.

To reproduce the archived run, execute
\begin{verbatim}
python verify.py --order 256 --independent-order 128
\end{verbatim}
The main exact checks require only the Python standard library. Installing
SymPy and mpmath enables the additional checks; missing optional packages
produce explicit skip messages. The files \texttt{verification\_output.txt}
and \texttt{verification\_metadata.json} record what actually ran.

\section{Connection with ProveIt and a formalization plan}
\label{bgg:sec:proveit}

The repository inspection was pinned to commit \repoSHA. The canonical
transseries-and-inversion README describes formal coefficient uniqueness,
quadratic/Catalan coefficients, remainder transport, and the distinction
between algebraic expansions and flat corrections \cite{proveit-volume}.
The source file \texttt{QuadraticCoreCatalan.lean} was also inspected
\cite{proveit-catalan}. It contains the Catalan description of
$\binom{1/2}{m}$ and finite convolution identities. It explicitly distinguishes
those finite algebraic results from a complete formal-power-series
existence/uniqueness package.

These are useful building blocks, not certificates for this article. In
particular, no claim here relies on a repository statement about A321941,
and no Lean or Rocq build was run for these new theorems.

A sensible formalization separates three layers. The first is finite
algebra: prove the Wronskian identity over a commutative ring and the
half-binomial weight formulas over $\Q$. The second is formal arithmetic:
define the coefficient arrays over $\Z[t]$, prove evenness, establish
uniqueness after coercion to $\Q[t]$, and formalize the modulo-$64$ argument
that descends to modulo $32$. The third is the analytic bridge: connect
actual Kummer/exponential-integral coefficients to the formal product
series by a full Poincar\'e expansion theorem.

The arithmetic layer can be formalized before the analytic one. Its theorem
would say that the uniquely defined formal solution has integral even
coefficients and the stated congruence. It should not be advertised as the
full analytic A321941 theorem until the expansion-identification bridge is
also verified. This separation makes the remaining trusted assumptions
visible rather than hiding them inside an informal numerical definition.

For implementation, \eqref{bgg:eq:algorithm-H}--\eqref{bgg:eq:algorithm-T} provide
finite sums only. The first proof can use a parity invariant asserting that
all stored coefficients except the constant term are even. The congruence
proof needs a second invariant tracking equality modulo $32$ while comparing
squares modulo $64$. The inspected Catalan module supports the alternative
fixed-point construction in \eqref{bgg:eq:catalan-square-root}. The inverse
arithmetic theorem requires only unit inversion, monic compositional
inversion, and the recursion $F^3=R^2$ over $\Z[1/3]$.

\section{Further research questions and topics}
\label{bgg:sec:future}

The following questions are not asserted as solved. They are arranged from
immediate extensions of the present arithmetic to broader analytic and
structural problems.

\begin{enumerate}[label=\textbf{\arabic*.}]
\item \textbf{Prove the remaining negativity conjecture.}
Establish $P_k(1)<0$ for all $k\ge3$. The exact formula for $P_k(0)$ gives a
negative reference edge, but mixed signs in the interior prevent a direct
coefficientwise argument. A positive integral representation for $-r_k$, a
sign-preserving reformulation, or an effective comparison with the constant
edge would resolve the missing assertion in the published Remark~11.

\emph{[Write note, batch 86A, 4 October 2026: answered.]} Part~\ref{bgg:sg:part:sign}
proves $r_k<0$ for every $k\ge3$ (Theorem~\ref{bgg:sg:thm:sign}),
not through the polynomials $P_k$ but through an exact radial-oscillator
heat kernel: $r_k$ has the sign of a coefficient $h_k$ dominated by the
single negative term $-\gamma_k$ of \eqref{bgg:sg:eq:hkdecomp} for
$k\ge32$, and an exact rational
certificate covers $3\le k\le31$. The dominating term is the constant
edge, $64^k(\tfrac12)_k\gamma_k=-P_k(0)$, so the proof is an effective
comparison with the constant edge in the sense suggested here.
The answer is AI-assisted and unrefereed.

\item \textbf{Determine all exact $2$-adic valuations.}
Theorem~\ref{bgg:thm:denominators-summary} handles indices with at most four
binary ones. Determine $v_2(r_k)$ at the remaining indices, including whether
there is a rule in terms of binary carries and a finite amount of additional
state. The existing proof gives a route through successively higher
moduli, but no such complete rule has been established here.

\item \textbf{Lift the universal congruence to higher powers of two.}
Find an effective description of $R_t$ modulo $2^m$ for arbitrary fixed $m$.
The valuations in \eqref{bgg:eq:weight-valuations} imply that only finitely many
shift distances are visible modulo each fixed power of two. Can this be
organized into a finite differential system, a recursive correction to an
algebraic series, or an efficient modular coefficient algorithm with proved
complexity?

\item \textbf{Study large-order growth of the coefficients themselves.}
The expansion parameter is $n\to\infty$, but the coefficient index
$k\to\infty$ is a different asymptotic problem. Determine a sharp growth law
for $r_k$ and the location and type of singularities of a suitable Borel
transform. The exact edge \eqref{bgg:eq:constant-polynomial} suggests a useful
comparison scale, but does not by itself prove the growth of $P_k(1)$.
Such a result would guide optimal truncation and could help prove eventual
negativity, leaving only finitely many indices.

\emph{[Write note, batch 86A, 4 October 2026: answered in part.]}
Part~\ref{bgg:sg:part:sign} proves the sharp growth law
$r_k\sim-16^k(k!)^2/(\pi^{3/2}k^{3/2})$, with a complete expansion in
inverse powers of $k$ relative to the edge $P_k(0)$
(Theorem~\ref{bgg:sg:thm:largeorder}), and it was used to prove
negativity outright. For the order-two Borel transform
$\sum_kd_ks^k/(2k)!$ it proves radius~$16$ and a logarithmic singularity
at $s=16$, as a radial boundary law (Theorem~\ref{bgg:sg:thm:borel}).
The analytic continuation, the full singularity set and Stokes data
stay open (Research question~\ref{bgg:sg:q:continuation}). The answer is
AI-assisted and unrefereed.

\item \textbf{Give effective analytic remainders and inverse certification.}
Combine the arithmetic recurrences with explicit bounds for the product of
Kummer functions. Determine an error bound after a specified truncation that
is valid for all $n$ above a computable threshold. Transport it through the
inverse map to obtain certified integer reconstruction or rounding. The
current all-orders theorem and high-precision checks do not supply these
constants.

\item \textbf{Analyze the parameter polynomials.}
Investigate the real zeros, sign intervals, coefficient triangles, and
large-$k$ asymptotics of $P_k(t)$. Determine for which fixed real or integer
$t$ the signs eventually stabilize. Establish analytic realizations of
\eqref{bgg:eq:lambda-coefficients} with uniform control in $\lambda$ away from
zero, and distinguish those from any singular scaling limit as
$\lambda\to0$.

\item \textbf{Extend the arithmetic mechanism to other recurrences.}
For normalized equations with $q_n=2+\lambda/n+\mu/n^2+\cdots$, classify the
potentials and normalizations for which the diagonal-product equation admits
integral coefficient operators after a single exponential scaling. The
quadratic invariant survives unchanged; the challenge is the arithmetic of
the shifted normalization and the resulting triangular equation. This is a
systematic source of further OEIS-related integrality problems rather than a
collection of isolated coefficient guesses.

\item \textbf{Find a signed combinatorial interpretation.}
The Catalan fixed point expresses the coefficients using trees built from
integer-weighted operations, but this is not yet a natural combinatorial
model for A321941. A model with a sign-reversing involution could explain both
dyadic scaling and the conjectured negative sign. It should account for the
parameter $t$ and the central-binomial reduction, not merely encode an
arbitrary integer recurrence tautologically.

\item \textbf{Complete a machine-checked proof and an independently reviewed edition.}
Formalize the finite arithmetic layer first, then the formal-analytic bridge.
An independent mathematical review should focus on the normalized Wronskian,
the cancellation leading to \eqref{bgg:eq:D-master}, the exact valuations of
shift weights, and the induction from modulus $64$ to modulus $32$. A
reviewed proof can then support a precise OEIS update distinguishing the
settled arithmetic assertions from the remaining sign question.
\end{enumerate}

\section{Conclusion}

The conjectural integrality of A321941 is not an accidental cancellation
visible only in a long list of coefficients. It follows from an exact
quadratic identity, integral rescaled shift operators, and a parity
induction. The same identity reduces modulo $64$ to a simple equation that
proves the central-binomial congruence modulo $32$.

The resulting structure is stronger than the original sequence statement:
it gives even integral parameter polynomials, explicit extreme coefficients,
minimal exponential scaling, infinitely many sharp denominator instances,
and an arithmetic restriction on an all-orders inverse expansion. The proof
is independent of finite verification ranges, while the accompanying checks
provide reproducible tests of every normalization used in the argument.

\appendix
\section{A short dependency and status ledger}

\begin{longtable}{@{}p{0.25\textwidth}p{0.40\textwidth}p{0.29\textwidth}@{}}
\toprule
Statement & Dependency & Status in this article\\
\midrule
\endhead
Original all-orders product expansion & Brent--Glasser--Guttmann,
Corollary 9 & Published input; cited, not re-proved\\
Quadratic identity & Elementary polynomial algebra & Fully proved here;
symbolic finite identity check included\\
Integral even $r_k$ & Formal bridge, weight arithmetic, triangular recurrence &
Fully proved here for every index\\
Congruence modulo $32$ & Evenness, operator valuations, reduction modulo $64$ &
Fully proved here for every index\\
Polynomial deformation & Parameterized formal equation & Fully proved here;
analytic identification only where an expansion is established\\
Denominator claims & Integrality, mod-$32$ result, factorial valuations &
Fully proved here within the stated binary-digit range\\
Inverse expansion & Published product expansion and formal reversion &
All fixed algebraic orders proved; no effective remainder constants\\
Negativity for $k\ge3$ & Would require additional sign arguments &
Not proved; finite observations only\\
Exponentially small sectors & Would require further analytic information &
Not determined\\
Lean/Rocq certification & Would require new formal developments and builds &
Not performed; a proposed route is described\\
\bottomrule
\end{longtable}

\emph{[Write note, batch 86A, 4 October 2026.]} The ledger describes
Part~\ref{bgg:part:arithmetic}. Negativity for $k\ge3$ is proved in
Part~\ref{bgg:sg:part:sign} (Theorem~\ref{bgg:sg:thm:sign}); exponentially
small sectors and Lean/Rocq certification remain as stated.

\section{Source audit and reproducibility scope}
\label{bgg:app:audit}

The source audit inspected the final journal PDF of \cite{BGG}, including
the conjecture and congruence on its printed page 10; the live A321941 entry;
the repository's transseries README; and the Catalan source module at the
pinned commit. Web searches included the exact sequence identifier with
``proof'', ``integrality'', and ``congruence'', as well as the authors and
paper title. This is a reproducible targeted audit, not an exhaustive search
of all publications, preprints, communications, or repository history.

The package contains the editable \LaTeX\ source, the compiled PDF, a
standalone verification script, the actual check log, coefficient data,
parameter-polynomial data, numerical diagnostics, and source notes. No
third-party paper or font file is redistributed. The script writes generated
outputs to its own directory by default; use \texttt{--out} to select another
location. Its main exact arithmetic does not require network access.

\emph{[Write note, batch 73O1.]} In the collection the compiled PDF is a
build of this text, not the delivered one; the script's own directory is
\texttt{code/}, so a run without \texttt{--out} would write its files
there, beside the program, and leave the records in \texttt{data/}
unchanged. The README gives a rerun on a copy.

\clearpage
\part{Signs, factorial divergence, and a Borel singularity}
\label{bgg:sg:part:sign}

% Part II continues the section numbers of Part I after its appendices.
\setcounter{section}{12}
\renewcommand{\thesection}{\arabic{section}}
\renewcommand{\theHsection}{sg.\arabic{section}}
\makeatletter\gdef\Hy@chapapp{\Hy@chapterstring}\makeatother

\section{The sign question and the large-order results}
\label{bgg:sg:sec:scope}

\emph{[Write note, batch 86A, 4 October 2026.]} This Part answers
question~1 of Section~\ref{bgg:sec:future} (the negativity $r_k<0$ for
every $k\ge3$) and the growth-law half of question~4 there. It is
printed from a later manuscript by the same pipeline; its provenance,
status and notation are in Section~\ref{bgg:sg:sec:collection}. Its
first section reproduces the A321941 portions of that manuscript's
opening section, which it shares with the manuscript's part on OEIS
A088714.

The manuscript's opening reads as follows. It develops two lines of
research suggested by the \emph{ProveIt} repository. The first concerns
the signs and large-order growth of the asymptotic coefficients
underlying OEIS A321941. The second concerns the much finer growth of
OEIS A088714, whose defining equation involves composition of its own
generating series; that line is printed as Part~V of the report
\path{a088714-bell-scale-growth}, not here. The organizing principle is
to replace a difficult global coefficient problem by an exact positive
representation, and then prove the uniform estimates that the
representation requires.

The repository was inspected at the fixed commit
\begin{center}\small\ttfamily
b7e4f25b6e79cc669fb245cba520017a0f4228b7.
\end{center}
The two reports used are Part~\ref{bgg:part:arithmetic} of the present
report and the A088714 report; their paths and Git blob identifiers are
recorded in Appendix~\ref{bgg:sg:app:sources}. Fixing the snapshot
matters: some conjectures in the original publications have already
been resolved in these reports, whereas others are explicitly left
open.

\subsection{Place in the ProveIt collection}
\label{bgg:sg:sec:collection}

\emph{[Write note, batch 86A, 4 October 2026.]}

\emph{Provenance.} This Part is built from one manuscript: Part~I
(Sections~2--7 and Appendix~A) of manuscript~02 of batch~86 of
ProveIt's incoming-reports intake, \emph{Signs, Factorial Divergence,
and Bell Normalization. New results for OEIS A321941, A088714, and
A088713}, dated 4 October 2026, author line ``ChatGPT'', in the archive
\texttt{oeis\_research\_bundle.zip}. The archive arrived in commit
\texttt{ae9baa422} and was placed by commit \texttt{0f084afa9}; its pin is
\texttt{b7e4f25b6}, at which this report's \texttt{article.tex} (blob
\texttt{10873a84}) was the text now printed as Part~\ref{bgg:part:arithmetic}.
The manuscript's Part~II (its Sections~8--12, the Bell normalization of
A088714 and A088713) is printed as Part~V of
\path{generating-functions-and-asymptotics/oeis-sequence-asymptotics/a088714-bell-scale-growth}.
Its shared material is divided by subject. Here are printed the A321941
portions of its Section~1 (Section~\ref{bgg:sg:sec:scope}), its
Sections~2--7 (Sections~\ref{bgg:sg:sec:resolvent}--\ref{bgg:sg:sec:borel}),
the A321941 rows and paragraphs of its Sections~13--14
(Section~\ref{bgg:sg:sec:changes}), its research questions~15.6--15.9
(Section~\ref{bgg:sg:sec:future}), its Appendix~A
(Appendix~\ref{bgg:sg:app:certificate}) and the A321941 parts of its
Appendix~B (Appendix~\ref{bgg:sg:app:sources}). Nothing was merged:
no other manuscript treats A321941. The manuscript's theorem
numbering is replaced by this report's: its Theorems~1.1, 1.2 and 7.1
are Theorems~\ref{bgg:sg:thm:sign}, \ref{bgg:sg:thm:largeorder}
and~\ref{bgg:sg:thm:borel}, and its Sections~2--7 are
Sections~14--19.

\emph{Edits.} The text is printed as delivered, except that every label
carries the prefix \texttt{bgg:sg:}, the symbols listed in
Table~\ref{bgg:sg:tab:notation} as renamed were renamed, references to
the manuscript's other parts, to its source report (now
Part~\ref{bgg:part:arithmetic}) and to its delivery files were
re-pointed, and notes marked \emph{[Write note]} were added. A
manuscript reference to the BGG article now cites the same
bibliography entry as Part~\ref{bgg:part:arithmetic}; the manuscript
accessed the OEIS entry A321941 on 4 October 2026.

\emph{Status.} The manuscript is AI-assisted and unrefereed, and nothing
in it is formalized in Lean, Rocq or any other proof assistant. Its
source audit records internal reviews (``independent analytic
normalization, cutoff32, all-orders and Borel reviews passed''), with
\texttt{external\_peer\_review}, \texttt{proof\_assistant\_formalized} and
\texttt{worldwide\_priority\_exhaustively\_verified} all false. Its
delivery README states that the infinite sign assertions, asymptotic
tails, convergence and error rates are proved in the manuscript and are
not inferred from finite tests, that no exhaustive worldwide priority
claim is made, and that the deliverable modifies neither the repository
nor the OEIS entries. The intake reran the five A321941 programs on a
copy; every recorded output was reproduced (up to Windows line
endings). Independently of the manuscript's code, it compared the
manuscript's heat coefficients with Part~\ref{bgg:part:arithmetic}'s own
coefficient table, computed by a different program from the quadratic
master equation: $d_k=(\tfrac12)_kh_k$ holds exactly for
$0\le k\le128$. From that table it also recomputed the middle column
and the last column of Table~\ref{bgg:sg:tab:signnumeric}, confirmed
$64^k(\tfrac12)_k\gamma_k=S_k$ for $1\le k<60$, $r_k<0$ for
$3\le k\le256$, and the exact margin $61/105$ of
\eqref{bgg:sg:eq:finitemargin}, attained at $k=3$; and it checked the
rational constants of Lemma~\ref{bgg:sg:lem:majorants} and of
Section~\ref{bgg:sg:sec:signproof} by hand. It did not re-derive the
special-function inputs of Sections~\ref{bgg:sg:sec:resolvent}
and~\ref{bgg:sg:sec:heat}.

\emph{Shipped layout.} The five programs of this Part are shipped as
\path{code/02-sign-*.py} and their recorded outputs as
\path{data/02-sign-*}, with the pinned SymPy requirement as
\path{data/02-sign-requirements-symbolic.txt}. The manuscript's driver,
which also runs the A088714 programs, its two run records and its
source audit are shipped once, with Part~V of the A088714 report. The
paths printed below are the delivered ones; the README gives reruns on
copies. The manuscript's \LaTeX\ source and PDF are not shipped and
survive in the arrival commit.

\emph{Relation to Part~\ref{bgg:part:arithmetic}.} The sign proof does
not use Part~\ref{bgg:part:arithmetic}'s integrality theorem. Its
comparison scale $S_k$ is exactly $-P_k(0)$ of
Theorem~\ref{bgg:thm:constant-edge}, the constant edge of the polynomial
deformation, and Theorem~\ref{bgg:sg:thm:largeorder} shows that this
edge is the leading asymptotic scale of $r_k=P_k(1)$:
$P_k(1)/P_k(0)\to1$. The Borel transform of
Section~\ref{bgg:sg:sec:borel} is radial information inside the disk of
convergence; continuation and the full singularity set remain open
(Section~\ref{bgg:sg:sec:future}), so question~4 of
Section~\ref{bgg:sec:future} is answered in its growth-law part only.

\emph{Notation.} Part~\ref{bgg:part:arithmetic} and this Part share
$a_n$, $b_n$, $\rho_n$, $d_k$, $r_k$, $k$, $n$, and Kummer's functions $M$
and $U$, with the same meanings. Every other symbol of this Part has
the meaning of Table~\ref{bgg:sg:tab:notation}, whose last column gives
the tempting false reading.

\begin{table}[ht]
\centering\small
\caption{Symbols of Part~II that also occur in
Part~\ref{bgg:part:arithmetic}. Renamed symbols are marked
\emph{(ms.\ $\cdot$)} with the manuscript's letter.}
\label{bgg:sg:tab:notation}
\begin{tabular}{@{}>{\raggedright\arraybackslash}p{0.17\textwidth}
>{\raggedright\arraybackslash}p{0.39\textwidth}
>{\raggedright\arraybackslash}p{0.36\textwidth}@{}}
\toprule
Symbol & Meaning in Part~II & \emph{Not}: meaning in Part~I\\
\midrule
$a_n$, $b_n$ & $[z^n]f_0$, $[z^n]f_1$ \emph{(ms.\ $\alpha_n$, $\eta_n$)} & --- (same objects)\\
$t$ & heat-kernel time & the deformation parameter of $P_k(t)$, used here only through $P_k(0)$, $P_k(1)$\\
$H(t)$, $h_k$ & the formal series \eqref{bgg:sg:eq:H} and its coefficients; $d_k=(\tfrac12)_kh_k$ & the operator $\Hop$; powers of the formal variable $h=1/n$ of \eqref{bgg:eq:D-definition}\\
$\gamma_j$ & comparison weights \eqref{bgg:sg:eq:qAbeta} \emph{(ms.\ $\beta_j$)} & the inverse coefficients $\beta_k$ of Theorem~\ref{bgg:thm:inverse}\\
$q(t)$ & $2\sinh(t/2)/t$ & $q_n=2+1/n$\\
$S_k$ & the scale \eqref{bgg:sg:eq:Sk}, $=-P_k(0)$ & the series $S(w)$, $S(z)$ in the proofs of Theorems~\ref{bgg:thm:polynomial-integrality}--\ref{bgg:thm:constant-edge}\\
$\mathcal M(t)$, $\mathcal M_\ell$ & coefficient majorant of Lemma~\ref{bgg:sg:lem:majorants} \emph{(ms.\ $M$)} & Kummer's $M(a,b,z)$, which keeps its letter\\
$N$ & fixed order in Theorem~\ref{bgg:sg:thm:largeorder} and Proposition~\ref{bgg:sg:prop:allorders} \emph{(ms.\ $R$)} & the series $R(z)$ of \eqref{bgg:eq:R-definition}\\
$\varepsilon_k$ & Borel remainder, proof of Theorem~\ref{bgg:sg:thm:borel} \emph{(ms.\ $e_k$)} & the weights $e_{j,\ell}$\\
$G_n(r,s)$, $K(t;2,2)$ & Green kernel, heat kernel & $G=eE_1(1)$\\
$F(n)$, $F_m(j)$ & $-2n^{3/2}\rho_n$; $[t^m]\mathcal A(t)q(t)^j$ & a generic series $F$; $F=R^{2/3}$\\
$W_r$, $W_z$ & Wronskians of the continuous solutions & the discrete Wronskian $W_n$\\
$\psi_m$, $\lambda_m$ & eigenfunctions, eigenvalues of $\mathscr L$ & $\psi(h)=D(h)^{-2/3}$; the parameter $\lambda$\\
$u$, $x$, $J$ & $16/k$, $1/k$, the split index & a solution $u_n$, $x_n=na_n$, the polynomial $J(h)$\\
$B$, $T$, $E$ (App.~\ref{bgg:sg:app:certificate}) & $q^{-1/2}$, $\tanh(t/4)$, $e^{-T}$ & $B(h)$, the array $T_n$, the operator $E(D)$\\
$c_j$, $C_{\!B}$, $\mathscr B_2$ & large-order coefficients, Borel constant, Borel transform & ---\\
\bottomrule
\end{tabular}
\end{table}

The regular solution $u_n(r)$ of Section~\ref{bgg:sg:sec:resolvent} and the
scalar $u=16/k$ of Section~\ref{bgg:sg:sec:signproof} share a letter in
the manuscript; the subscript and argument distinguish them.

\subsection{The A321941 question}

Let
\[
 f_0(z)=\exp\!\left(\frac{z}{1-z}\right),\qquad
 f_1(z)=\exp\!\left(\frac1{1-z}\right)
                E_1\!\left(\frac1{1-z}\right),
\]
where \(E_1(x)=\int_x^\infty e^{-u}u^{-1}\,du\), and put
\(a_n=[z^n]f_0(z)\), \(b_n=[z^n]f_1(z)\), and
\(\rho_n=a_nb_n\).
Brent, Glasser, and Guttmann \cite{BGG}, abbreviated BGG, proved the
Poincare expansion
\begin{equation}\label{bgg:sg:eq:rho}
 \rho_n\sim-\frac1{2n^{3/2}}\sum_{k\ge0}d_kn^{-k},
 \qquad d_0=1.
\end{equation}
Set \(r_k=64^kd_k\).  Strictly speaking, A321941 lists the
\emph{numerators} of the rational numbers \(r_k\)
\cite{oeis321941}.  The integrality theorem of Part~\ref{bgg:part:arithmetic}
(Theorem~\ref{bgg:thm:main}) identifies
these numerators with \(r_k\) itself.  The sign argument here does not
require that integrality theorem: a rational number and its numerator
have the same sign.

BGG's Remark~11 reports negativity for \(3\le k\le1000\).
The repository report (Part~\ref{bgg:part:arithmetic} here) proves
integrality and congruences, but explicitly
leaves negativity at every later index open.  We prove:

\begin{theorem}[Complete sign pattern]\label{bgg:sg:thm:sign}
The coefficients in \eqref{bgg:sg:eq:rho} satisfy
\[
 r_0=1,\qquad r_1=-14,\qquad r_2=86,\qquad
 \boxed{r_k<0\quad\hbox{for every integer }k\ge3.}
\]
\end{theorem}

The proof supplies more than eventual negativity.  It gives a uniform
inequality for every \(k\ge32\), with the remaining \(29\) cases
certified in exact rational arithmetic.  An exact spectral
identification, proved below, establishes that the coefficients being
bounded are precisely those of \eqref{bgg:sg:eq:rho}.

Define the positive comparison scale
\begin{equation}\label{bgg:sg:eq:Sk}
 S_k=\binom{2k}{k}(2k-3)!!(2k+1)!!\qquad(k\ge1),
\end{equation}
with the convention \((-1)!!=1\).
Our second result describes the expansion \emph{index} \(k\) tending
to infinity, rather than the original variable \(n\).

\begin{theorem}[Large-order expansion]\label{bgg:sg:thm:largeorder}
There are effectively computable rational numbers \(c_j\), with
\(c_0=1\), such that, for each fixed \(N\ge0\),
\[
 r_k=-S_k\left(\sum_{j=0}^{N}\frac{c_j}{k^j}
                       +O_N(k^{-N-1})\right).
\]
The first terms are
\begin{equation}\label{bgg:sg:eq:largeorder4}
 \frac{r_k}{-S_k}
 =1-\frac1{3k}-\frac{11}{18k^2}
   -\frac{1279}{1620k^3}-\frac{5089}{9720k^4}
   +O(k^{-5}).
\end{equation}
In particular,
\[
 r_k\sim-\frac{16^k(k!)^2}{\pi^{3/2}k^{3/2}},
 \qquad
 d_k\sim-\frac{(k!)^2}{\pi^{3/2}4^kk^{3/2}}.
\]
\end{theorem}

These estimates determine the exact order of factorial divergence and
the first positive singularity of an explicitly normalized Borel
transform; see Theorem~\ref{bgg:sg:thm:borel}.

\subsection{What the proofs and computations establish}
\label{bgg:sg:sec:establish}

The mathematical proofs are given in the manuscript.  The accompanying
programs check exact finite identities, the finite sign certificate,
and independently generated asymptotic coefficients.  Their role is
specific: no finite experiment substitutes for a uniform estimate over
an infinite tail.

The conclusions are advances beyond the inspected reports and the
primary literature cited here.  This statement concerns that source
boundary; it is not an exhaustive claim of worldwide priority.
The manuscript is a research proof with independent checks, rather
than a proof-assistant formalization or a claim of external peer review.
The repository's existing arithmetic theorems are credited as prior
results throughout.

\section{An exact radial-oscillator representation}\label{bgg:sg:sec:resolvent}

The key step is an exact representation of the product in
\eqref{bgg:sg:eq:rho}.  BGG's Lemmas~1--2 give, for \(n\ge1\),
\begin{equation}\label{bgg:sg:eq:kummercoeff}
 a_n=e^{-1}M(n+1,2,1),\qquad
 b_n=-\Gamma(n)U(n,0,1),
\end{equation}
where \(M\) and \(U\) are Kummer functions \cite{BGG,dlmf-kummer}.

Consider the Friedrichs realization on \(L^2((0,\infty),dr)\) of
\begin{equation}\label{bgg:sg:eq:operator}
 \mathscr L=-\frac{d^2}{dr^2}+\frac{r^2}{16}+\frac{3}{4r^2}.
\end{equation}
The regular solution at zero and recessive solution at infinity
of \((\mathscr L+n)y=0\) are, respectively,
\[
 u_n(r)=r^{3/2}e^{-r^2/8}M(n+1,2,r^2/4),\qquad
 v_n(r)=r^{3/2}e^{-r^2/8}U(n+1,2,r^2/4).
\]
These statements follow by substituting \(z=r^2/4\) in Kummer's
equation.  The behaviors at the two endpoints also specify the
Green kernel directly, if one prefers to start from the differential
equation.

The Wronskian identity and transformation
\[
 W_z\{M(n+1,2,z),U(n+1,2,z)\}
       =-\frac{e^zz^{-2}}{\Gamma(n+1)},\qquad
 U(n,0,z)=zU(n+1,2,z)
\]
are DLMF~13.2.34 and 13.2.40 \cite{dlmf-kummer}.  The first
formula uses ordinary \(M\); the regularized convention in DLMF agrees
here because \(\Gamma(2)=1\).
Multiplication by the gauge factor squared and \(dz/dr=r/2\)
gives
\[
 W_r\{u_n,v_n\}=-\frac8{\Gamma(n+1)}.
\]
Consequently the kernel of \((\mathscr L+n)^{-1}\) is
\[
 G_n(r,s)=\frac{\Gamma(n+1)}8u_n(\min(r,s))v_n(\max(r,s)).
\]
At \(r=s=2\), this becomes
\begin{equation}\label{bgg:sg:eq:greenrho}
 G_n(2,2)=e^{-1}\Gamma(n+1)M(n+1,2,1)U(n,0,1)
         =-n\rho_n.
\end{equation}
Thus the exact normalization needed for \eqref{bgg:sg:eq:rho} is
\begin{equation}\label{bgg:sg:eq:Fgreen}
 F(n):=-2n^{3/2}\rho_n=2\sqrt n\,G_n(2,2).
\end{equation}
In particular, every constant in the coefficient identification is
fixed before any asymptotic expansion is made.

\section{The heat kernel and its formal coefficients}\label{bgg:sg:sec:heat}

\subsection{Summing the spectral kernel}

The normalized eigenfunctions and eigenvalues of \(\mathscr L\) are
\[
 \psi_m(r)=\frac{r^{3/2}e^{-r^2/8}}{\sqrt{8(m+1)}}
                     L_m^{(1)}(r^2/4),\qquad \lambda_m=m+1.
\]
The Laguerre equation verifies the eigenvalue.  Their normalization
follows from
\[
 \int_0^\infty r^3e^{-r^2/4}L_m^{(1)}(r^2/4)^2\,dr
 =8\int_0^\infty xe^{-x}L_m^{(1)}(x)^2\,dx=8(m+1).
\]
Laguerre completeness gives a basis \cite{dlmf-laguerre}.  Hence
\[
 K(t;2,2)
 =e^{-1}\sum_{m\ge0}\frac{L_m^{(1)}(1)^2}{m+1}e^{-(m+1)t}.
\]
For \(0<z<1\), the Hille--Hardy identity is
\[
 \sum_{m\ge0}\frac{L_m^{(1)}(1)^2}{m+1}z^m
 =\frac{z^{-1/2}}{1-z}\exp\!\left(-\frac{2z}{1-z}\right)
        I_1\!\left(\frac{2\sqrt z}{1-z}\right);
\]
this is DLMF~18.18.27 with both arguments equal to one
\cite{dlmf-laguerre}.  Substitution of \(z=e^{-t}\) yields the exact
formula
\begin{equation}\label{bgg:sg:eq:exactheat}
 \boxed{
 K(t;2,2)=
 \frac{\exp[-\coth(t/2)]}{2\sinh(t/2)}
            I_1(\operatorname{csch}(t/2)),\qquad t>0.}
\end{equation}
The nonnegative spectral summands allow Tonelli's theorem to give
\begin{equation}\label{bgg:sg:eq:laplace}
 G_n(2,2)=\int_0^\infty e^{-nt}K(t;2,2)\,dt.
\end{equation}
The heat kernel is \(O(t^{-1/2})\) near zero and decays exponentially
at infinity, so this integral is finite for \(n>0\).

\subsection{A coefficient formula with a negative leading contribution}

Introduce
\begin{equation}\label{bgg:sg:eq:qAbeta}
 q(t)=\frac{2\sinh(t/2)}t,\qquad
 \mathcal A(t)=q(t)^{-1/2}e^{-\tanh(t/4)},\qquad
 \gamma_j=\frac{(2j-3)!!(2j+1)!!}{16^j j!}\quad(j\ge1).
\end{equation}
The calligraphic symbol \(\mathcal A\) belongs only to this Part; it is
different from the A088714 series \(A\) of the manuscript's other part.
The large-positive-argument expansion of \(I_1\) is
\[
 I_1(x)\sim\frac{e^x}{\sqrt{2\pi x}}
 \left(1-\sum_{j\ge1}
       \frac{(2j-3)!!(2j+1)!!}{8^jj!\,x^j}\right)
 \quad\cite{dlmf-bessel}.
\]
Every displayed correction is negative.
Using
\(\coth(t/2)-\operatorname{csch}(t/2)=\tanh(t/4)\),
define the formal series
\begin{equation}\label{bgg:sg:eq:H}
 H(t)=\mathcal A(t)
       \left(1-\sum_{j\ge1}\gamma_jt^jq(t)^j\right)
      =\sum_{k\ge0}h_kt^k.
\end{equation}
For each finite order the Bessel expansion in \eqref{bgg:sg:eq:exactheat}
proves
\[
 K(t;2,2)\sim(4\pi t)^{-1/2}\sum_{k\ge0}h_kt^k
 \qquad(t\downarrow0).
\]
Equation \eqref{bgg:sg:eq:H} is a \emph{formal} identity.  It does not assert
convergence of the infinite Bessel expansion or of \(H\).

Integrating a finite truncation in \eqref{bgg:sg:eq:laplace} and bounding its
remainder gives Watson's lemma \cite{dlmf-watson}.  In combination
with \eqref{bgg:sg:eq:Fgreen}, it proves
\[
 F(n)\sim\sum_{k\ge0}\frac{\Gamma(k+\tfrac12)}{\sqrt\pi}h_kn^{-k}.
\]
Uniqueness of Poincare coefficients now gives the exact relation
\begin{equation}\label{bgg:sg:eq:dh}
 \boxed{d_k=(\tfrac12)_k h_k.}
\end{equation}
In particular, the sign of \(r_k\) is the sign of \(h_k\).
The first values, directly from \eqref{bgg:sg:eq:H}, are
\begin{equation}\label{bgg:sg:eq:hinitial}
 h_0=1,\qquad h_1=-\frac7{16},\qquad
 h_2=\frac{43}{1536},\qquad h_3=-\frac{61}{8192}.
\end{equation}

\section{Coefficient majorants}\label{bgg:sg:sec:majorants}

For a formal series \(P\), write \(|P|\) for the series of absolute
values of its coefficients.  We write \(P\preccurlyeq Q\) when
coefficientwise domination holds and \(Q\) has nonnegative coefficients.

\begin{lemma}\label{bgg:sg:lem:majorants}
The following coefficient inequalities hold:
\[
 q(t)\preccurlyeq e^{t^2/24},\qquad
 |\mathcal A(t)|\preccurlyeq
 \mathcal M(t):=(2-q(t))^{-1/2}e^{\tan(t/4)}.
\]
The series \(\mathcal M\) has nonnegative coefficients, \(\mathcal M(0)=1\), and
\begin{equation}\label{bgg:sg:eq:Mone}
 \mathcal M(1)<\frac32.
\end{equation}
Moreover \(q(1)\le24/23\).
\end{lemma}

\begin{proof}
The coefficient of \(t^{2m}\) in \(q\) is
\(1/(4^m(2m+1)!)\).  Thus the first assertion is equivalent to
\(6^mm!\le(2m+1)!\), which follows by induction.

Write \(u=q-1\), so \(u\) has nonnegative coefficients.  The absolute
binomial coefficients of \((1+u)^{-1/2}\) are the coefficients of
\((1-u)^{-1/2}\).  Similarly, the absolute coefficients of
\(\tanh(t/4)\) are those of \(\tan(t/4)\).  Positivity of the
latter coefficients also follows recursively from
\(T'=(1+T^2)/4\), \(T(0)=0\), for \(T(t)=\tan(t/4)\).
Multiplication and exponentiation preserve coefficient majorants,
giving the asserted bound by \(M\).

Here are rational bounds for \eqref{bgg:sg:eq:Mone}.  The first nonconstant
term in \(q(1)\) is \(1/24\), and the ratio of successive terms
thereafter is at most \(1/80\).  Therefore
\[
 q(1)\le1+\frac{1/24}{1-1/80}=\frac{247}{237}<\frac{21}{20}.
\]
It follows that \((2-q(1))^{-1/2}<\sqrt{20/19}<11/10\).
The elementary estimates
\(\sin(1/4)\le1/4\) and \(\cos(1/4)\ge31/32\) give
\(\tan(1/4)\le8/31\).  Since \(e^x\le(1-x)^{-1}\)
for \(0\le x<1\),
\[
 \mathcal M(1)<\frac{11}{10}\frac{31}{23}=\frac{341}{230}<\frac32.
\]
Finally \(q(1)\le e^{1/24}\le24/23\).
\end{proof}

\begin{lemma}\label{bgg:sg:lem:beta}
The positive numbers \(\gamma_j\) satisfy
\begin{equation}\label{bgg:sg:eq:betaratio}
 \frac{\gamma_j}{\gamma_{j-1}}
  =\frac{(2j-3)(2j+1)}{16j}\ge\frac j8
  \qquad(j\ge3).
\end{equation}
They decrease through \(\gamma_5\) and increase thereafter.
For every \(J\ge15\), \(\gamma_J=\max_{1\le j\le J}\gamma_j\), and
\[
 \gamma_{16}=
 \frac{61394155884765868567528125}
      {604462909807314587353088}>100.
\]
\end{lemma}

\begin{proof}
The ratio follows from \eqref{bgg:sg:eq:qAbeta}; its lower bound is equivalent
to \(2j^2-4j-3\ge0\) for \(j\ge3\).  The same ratio proves the
monotonicity assertions.  The displayed rational evaluation gives
\(\gamma_{16}>100\), while
\(\gamma_{16}/\gamma_{15}=957/256<4\).  Hence
\(\gamma_{15}>25>\gamma_1=3/16\), which settles the maximum assertion.
\end{proof}

\section{Negativity at every index}\label{bgg:sg:sec:signproof}

\subsection{A uniform bound beyond index 31}

Extracting the coefficient of \(t^k\) in \eqref{bgg:sg:eq:H} gives
\begin{equation}\label{bgg:sg:eq:hkdecomp}
 h_k=-\gamma_k+\mathcal A_k
       -\sum_{j=1}^{k-1}\gamma_j[t^{k-j}]\mathcal A(t)q(t)^j,
 \qquad \mathcal A_k=[t^k]\mathcal A(t).
\end{equation}
We show that the absolute value of all terms other than
\(-\gamma_k\) is less than \(9\gamma_k/10\).

Fix \(k\ge32\), set \(u=16/k\le1/2\), and split the sum at
\(J=\lceil k/2\rceil-1\).
For the near indices \(j\ge\lceil k/2\rceil\), put \(m=k-j\).
The ratio bound \eqref{bgg:sg:eq:betaratio} gives
\(\gamma_{k-m}/\gamma_k\le u^m\).
Also
\[
 |\mathcal A q^j|\preccurlyeq \mathcal M(t)e^{kt^2/24}.
\]
Consequently their total normalized absolute contribution is at most
\begin{equation}\label{bgg:sg:eq:nearerror}
 \mathcal M(u)e^{ku^2/24}-1.
\end{equation}
Since \(M\) has nonnegative coefficients and \(u^\ell\le u\) for
\(\ell\ge1\),
\[
 \mathcal M(u)\le1+u(\mathcal M(1)-1)<1+\frac u2=1+\frac8k.
\]
Using \(ku^2/24=32/(3k)<1\) and \(e^x\le(1-x)^{-1}\),
\begin{equation}\label{bgg:sg:eq:nearbound}
 \mathcal M(u)e^{ku^2/24}-1
 <\frac{1+8/k}{1-32/(3k)}-1
 =\frac{56}{3k-32}\le\frac78.
\end{equation}

For the far indices \(1\le j\le J\), Lemma~\ref{bgg:sg:lem:beta} applies
because \(J\ge15\).  It gives
\[
 \frac{\gamma_j}{\gamma_k}\le
 \frac{\gamma_J}{\gamma_k}\le(16/k)^{k/2}.
\]
Evaluation of the nonnegative majorant at one gives
\[
 |[t^{k-j}]\mathcal A q^j|
 \le \mathcal M(1)q(1)^j\le\frac32(24/23)^{k/2}.
\]
There are fewer than \(k/2\) far indices, so their combined
normalized contribution is bounded by
\begin{equation}\label{bgg:sg:eq:farbound}
 \frac{3k}{4}\left(\frac{384}{23k}\right)^{k/2}
 <\frac1{100}\qquad(k\ge32).
\end{equation}
For completeness, the expression on the left decreases at successive
integers: its ratio is at most
\((33/32)\sqrt{12/23}<1\).  At \(k=32\) it equals
\(24(12/23)^{16}<1/100\), an exact rational inequality.

Finally
\[
 \frac{|\mathcal A_k|}{\gamma_k}<\frac{3}{200},
\]
by \eqref{bgg:sg:eq:Mone} and \(\gamma_k\ge\gamma_{16}>100\).
Combining the three contributions proves
\begin{equation}\label{bgg:sg:eq:strongsign}
 \left|\frac{h_k}{\gamma_k}+1\right|
 <\frac78+\frac1{100}+\frac3{200}=\frac9{10}
 \qquad(k\ge32).
\end{equation}
In particular
\[
 -\frac{19}{10}\gamma_k<h_k<-\frac1{10}\gamma_k<0.
\]

\subsection{The finite certificate}

All \(h_k\) in \eqref{bgg:sg:eq:H} are rational.  Appendix~\ref{bgg:sg:app:certificate}
gives a short executable certificate for \(3\le k\le31\),
using only integer and rational arithmetic.  The supplied full verifier
checks \(3\le k\le128\); a separately implemented recurrence from
BGG agrees with the heat construction through degree \(64\).
The exact minimum in the required finite range is
\begin{equation}\label{bgg:sg:eq:finitemargin}
 \min_{3\le k\le31}\frac{-h_k}{\gamma_k}=\frac{61}{105}>0.
\end{equation}
Thus every finite case is decided by an exact comparison.
Together with \eqref{bgg:sg:eq:strongsign}, \eqref{bgg:sg:eq:dh}, and
\eqref{bgg:sg:eq:hinitial}, this proves Theorem~\ref{bgg:sg:thm:sign}.

The initial scaled values provide a useful normalization check:
\[
\begin{split}
 r_0,\ldots,r_6={}&1,\ -14,\ 86,\ -3660,\ -1042202,\\
 &-247948260,\ -108448540420.
\end{split}
\]
They agree with the defining OEIS entry \cite{oeis321941}.

\section{The complete expansion as \texorpdfstring{\(k\)}{k} grows}\label{bgg:sg:sec:largeorderproof}

\subsection{A finite formula at each asymptotic order}

For fixed \(m\ge0\), define
\[
 F_m(j)=[t^m]\mathcal A(t)q(t)^j.
\]
Because \(\log q(t)\) starts at degree two, \(F_m(j)\) is a
polynomial in \(j\) of degree at most \(\lfloor m/2\rfloor\).
For fixed \(m\) and \(k>m\),
\begin{equation}\label{bgg:sg:eq:fixedmratio}
 \frac{\gamma_{k-m}}{\gamma_k}
 =\prod_{r=0}^{m-1}
       \frac{16(k-r)}{(2(k-r)-3)(2(k-r)+1)}.
\end{equation}
Its order is \(k^{-m}\).  Thus the product
\((\gamma_{k-m}/\gamma_k)F_m(k-m)\) starts at order
\(k^{-\lceil m/2\rceil}\).

\begin{proposition}\label{bgg:sg:prop:allorders}
For every fixed integer \(N\ge0\),
\begin{equation}\label{bgg:sg:eq:allordersfinite}
 \frac{h_k}{-\gamma_k}
 =\sum_{m=0}^{2N}
       \frac{\gamma_{k-m}}{\gamma_k}F_m(k-m)
       +O_N(k^{-N-1}).
\end{equation}
For \(m=0\), the summand is understood to be one.
\end{proposition}

\begin{proof}
Use the same far/near split as in
Section~\ref{bgg:sg:sec:signproof}.  The far bound
\eqref{bgg:sg:eq:farbound} is smaller than every power of \(k^{-1}\).
The pure analytic contribution also has that property: from
\eqref{bgg:sg:eq:betaratio},
\[
 \gamma_k\ge\frac{\gamma_2}{2}\,8^{2-k}k!,
 \qquad |\mathcal A_k|<\frac32.
\]

For the remaining near terms with \(m\ge2N+1\), expand
\(\mathcal M(t)=\sum_{\ell\ge0}\mathcal M_\ell t^\ell\).  Their majorant is
\[
 \sum_{\ell+2v\ge2N+1}
 \mathcal M_\ell u^\ell\frac{(ku^2/24)^v}{v!},
 \qquad u=\frac{16}{k}.
\]
In terms of \(x=1/k\), this is the nonnegative analytic series
\begin{equation}\label{bgg:sg:eq:constrainedmajorant}
 \sum_{\ell+2v\ge2N+1}
 \mathcal M_\ell16^\ell\left(\frac{32}{3}\right)^v
                  \frac{x^{\ell+v}}{v!}.
\end{equation}
Every exponent satisfies \(\ell+v\ge N+1\).
The unconstrained series is
\(\mathcal M(16x)e^{32x/3}\), analytic near zero.
Therefore \eqref{bgg:sg:eq:constrainedmajorant} is \(O_N(x^{N+1})\).
Keeping the finitely many remaining terms proves
\eqref{bgg:sg:eq:allordersfinite}.
\end{proof}

Expanding the rational functions in \eqref{bgg:sg:eq:fixedmratio}, and the
finite polynomials \(F_m\), gives an algorithm for every coefficient
in Theorem~\ref{bgg:sg:thm:largeorder}.  No divergent series is summed in
this calculation.  Through \(m=8\), it gives
\[
 \frac{h_k}{-\gamma_k}
 =1-\frac1{3k}-\frac{11}{18k^2}
  -\frac{1279}{1620k^3}-\frac{5089}{9720k^4}+O(k^{-5}).
\]
Two independent exact implementations reproduce these coefficients.
The primary generator additionally records
\[
 c_5=\frac{124787}{81648},\qquad
 c_6=\frac{368256067}{36741600},\qquad
 c_7=\frac{8769001777}{220449600},\qquad
 c_8=\frac{173295857159}{1322697600}.
\]

\subsection{Factorial scale and numerical diagnostics}

The exact identity
\begin{equation}\label{bgg:sg:eq:scalegamma}
 64^k(\tfrac12)_k\gamma_k
 =S_k
 =\frac{16^k}{\pi^{3/2}}
   \frac{\Gamma(k-\tfrac12)\Gamma(k+\tfrac12)
         \Gamma(k+\tfrac32)}{\Gamma(k+1)}
\end{equation}
converts Proposition~\ref{bgg:sg:prop:allorders} into
Theorem~\ref{bgg:sg:thm:largeorder}.  Stirling's formula gives
\[
 S_k\sim\frac{16^k(k!)^2}{\pi^{3/2}k^{3/2}},
 \qquad \frac{r_{k+1}}{r_k}\sim16k^2.
\]
Thus both ordinary formal series \(\sum d_kz^k\) and
\(\sum r_kz^k\) have radius zero and exact Gevrey order two:
their coefficients grow at the scale of a fixed exponential times
\((k!)^2\), and no lower factorial order can bound them.

Table~\ref{bgg:sg:tab:signnumeric} illustrates the proved expansion.  Each
entry in its middle column is rounded only after exact rational
evaluation.  The last column subtracts the displayed four-correction
approximation \eqref{bgg:sg:eq:largeorder4}.  These values check the scale
and normalization; the uniform remainder was proved above.

\begin{table}[htbp]
\centering
\caption{Large-order diagnostics for A321941.}\label{bgg:sg:tab:signnumeric}
\begin{tabular}{rrr}
\toprule
\(k\)&\(r_k/(-S_k)\)&Difference from \eqref{bgg:sg:eq:largeorder4}\\
\midrule
3   &0.580952380952&\(-2.0433\,10^{-1}\)\\
4   &1.050342151675&\( 1.8625\,10^{-1}\)\\
8   &0.947282773338&\( 1.6788\,10^{-4}\)\\
16  &0.976581013320&\( 2.2388\,10^{-6}\)\\
32  &0.988962008187&\( 5.6170\,10^{-8}\)\\
64  &0.994639428263&\( 1.5788\,10^{-9}\)\\
128 &0.997358155702&\( 4.6832\,10^{-11}\)\\
\bottomrule
\end{tabular}
\end{table}

\section{An order-two Borel singularity}\label{bgg:sg:sec:borel}

We fix the factorial convention explicitly:
\begin{equation}\label{bgg:sg:eq:BorelDef}
 \mathscr B_2(s)=\sum_{k\ge0}\frac{d_k}{(2k)!}s^k.
\end{equation}

\begin{theorem}\label{bgg:sg:thm:borel}
The power series \eqref{bgg:sg:eq:BorelDef} has radius of convergence exactly
\(16\).  There is a finite real constant \(C_{\!B}\) such that
\begin{equation}\label{bgg:sg:eq:borelboundary}
 \mathscr B_2(16x)
 =\frac1\pi\log(1-x)+C_{\!B}
    +O\bigl((1-x)|\log(1-x)|\bigr)
 \qquad(x\uparrow1).
\end{equation}
In particular the positive boundary point \(s=16\) is singular.
\end{theorem}

\begin{proof}
The gamma duplication formula and \eqref{bgg:sg:eq:scalegamma} give
\[
 \frac{(\tfrac12)_k\gamma_k}{(2k)!}
 =\frac{\Gamma(k-\tfrac12)\Gamma(k+\tfrac32)}
       {\pi16^k\Gamma(k+1)^2}.
\]
Theorem~\ref{bgg:sg:thm:largeorder} therefore yields
\[
 \frac{16^kd_k}{(2k)!}=-\frac1{\pi k}+O(k^{-2}).
\]
The root test proves the asserted radius.  Let
\[
 \varepsilon_k=\frac{16^kd_k}{(2k)!}+\frac1{\pi k},\qquad
 C_{\!B}=1+\sum_{k\ge1}\varepsilon_k.
\]
The constant is well defined because \(\sum|\varepsilon_k|<\infty\).
For \(0<x<1\), use
\(-\sum_{k\ge1}x^k/k=\log(1-x)\).  The remaining boundary error
is bounded by a constant times
\[
 \sum_{k\ge1}\frac{1-x^k}{k^2}
 =O\bigl((1-x)|\log(1-x)|\bigr).
\]
To verify the last estimate directly, split at
\(k=\lfloor(1-x)^{-1}\rfloor\) and use
\(1-x^k\le\min\{k(1-x),1\}\).
\end{proof}

This is a radial boundary statement proved inside the convergence
disk.  It makes no claim of analytic continuation or Borel
summability.  With the same divisor \((2k)!\), the transform of
\(r_k=64^kd_k\) has radius \(1/4\), rather than \(16\).
The change of radius is entirely due to the normalization.

\section{What has changed, and what the programs check}
\label{bgg:sg:sec:changes}

\emph{[Write note, batch 86A.]} This section prints the A321941
portions of the manuscript's Sections~13 and~14, which are shared with
its A088714 part. The three A088714 and A088713 rows of its table, its
paragraph on A088714, its A088714 check items and its normalization
diagnostics are printed in Part~V of the A088714 report.

\subsection{What has changed mathematically}

The source boundary and the new conclusions for A321941 are summarized
in Table~\ref{bgg:sg:tab:boundary}.  The strongest improvements concern
infinite signs and multiplicative normalization.  In the first problem,
integrality alone gives little control over sign.  In the second,
an \(n\)-th-root asymptotic leaves room for multiplicative factors
whose logarithms diverge.  The new proofs supply exactly the missing
global control in each case.

\begin{table}[htbp]
\centering\small
\caption{Earlier results and the conclusions established here (A321941
rows of the manuscript's table).}
\label{bgg:sg:tab:boundary}
\begin{tabularx}{\textwidth}{@{}>{\raggedright\arraybackslash}p{23mm}>{\raggedright\arraybackslash}X>{\raggedright\arraybackslash}X@{}}
\toprule
Topic&Inspected earlier boundary&This article\\
\midrule
A321941 signs&
BGG observed negativity through \(1000\); the repository leaves the
infinite assertion open.&
Every \(r_k<0\) for \(k\ge3\), with an explicit bound for all \(k\ge32\).\\[3pt]
A321941 growth&
The repository gives an exact constant-edge expression for its
polynomial deformation.&
That expression is the leading asymptotic scale, with a complete
inverse-\(k\) expansion and an order-two Borel boundary law.\\
\bottomrule
\end{tabularx}
\end{table}

For A321941 the spectral representation turns the sign question into
domination by a single explicit coefficient.  Near offsets determine
the entire large-order expansion, while far offsets are negligible
beyond every algebraic order.  The formal-series viewpoint is essential:
the divergent series is used coefficient by coefficient, after the
exact resolvent and heat kernel have already been established.

\subsection{What is checked exactly}

The primary verification programs use the Python standard library.
A separate optional symbolic cross-check uses SymPy.
The A321941 routines use \texttt{fractions.Fraction}; their sign
decisions never involve floating-point arithmetic.  The included
A321941 checks cover:
\begin{enumerate}
\item The required finite sign range \(3\le k\le31\), with an
extended check through \(128\).
\item Agreement with BGG's published coefficient recurrence through
degree \(64\), and agreement with all \(13\) initial terms displayed
in the inspected OEIS entry.
\item The rational constants in the cutoff-\(32\) proof, including
\(24(12/23)^{16}<1/100\) and the combined error \(9/10<1\).
\item Independent constructions of the large-order coefficients
through \(c_8\).
\end{enumerate}

The tests are finite certificates and checks of algebraic
identifications.  They do not establish the infinite asymptotic
remainder bounds; those were proved in
Sections~\ref{bgg:sg:sec:resolvent}--\ref{bgg:sg:sec:borel} (and, for
A088714, in Part~V of that report).
The independent A321941 implementation starts from BGG's polynomial
coefficient recurrence and reconstructs the heat coefficients by a
different route.  This checks more than agreement of two instances
of the same series-expansion code.

\subsection{Rebuilding the deliverable}

The archive contains a self-contained \texttt{oeis\_research.tex},
its compiled PDF, the verification code, exact coefficient data,
recorded outputs, and a source audit.  The TeX source has no
external figures or bibliography file.
Compile it with a standard TeX Live installation:
\begin{lstlisting}[language=bash]
latexmk -pdf -interaction=nonstopmode -halt-on-error oeis_research.tex
\end{lstlisting}
The file \texttt{README.txt} gives the individual verification commands.
Python~3.11 or later supports the Decimal methods used by the optional
normalization diagnostics.  The exact rational sign checks and
integer endpoint checks need no third-party numerical package.

\emph{[Write note, batch 86A.]} The delivered \texttt{oeis\_research.tex},
its PDF and its \texttt{README.txt} are not shipped; they survive in the
arrival commit \texttt{ae9baa422}. In the collection the A321941
programs are \path{code/02-sign-finite_certificate.py},
\path{02-sign-verify_sign.py}, \path{02-sign-independent_check.py},
\path{02-sign-large_order.py} and \path{02-sign-asymptotic_check.py},
and their recorded outputs are in \path{data/} under the same prefix.
Two of them find a sibling by its delivered name
(\path{large_order.py} imports \path{verify_sign},
\path{asymptotic_check.py} loads \path{independent_check.py}), four
write their output beside themselves, and the manuscript's driver
\path{run_checks.py}, shipped with Part~V of the A088714 report,
expects the delivered layout and rewrites its own run records. The
README of this report therefore gives reruns on copies with the
delivered names. The PDF of this report is a build of this text.

\section{Further research questions for Part~II}
\label{bgg:sg:sec:future}

\emph{[Write note, batch 86A.]} These are the manuscript's research
questions~15.6--15.9; its questions~15.1--15.5 concern A088714 and are
printed in Part~V of the A088714 report. Question~\ref{bgg:sg:q:region}
continues question~6 of Section~\ref{bgg:sec:future},
question~\ref{bgg:sg:q:continuation} is the part of question~4 there
that this Part leaves open, and question~\ref{bgg:sg:q:arithmetic}
bears on questions~2 and~3 there.

The following questions are deliberately more precise than a request
for additional numerical terms.  Each identifies a mathematical
target that is not established by the present proofs.

\begin{question}[The entire parameter-dependent sign region]
\label{bgg:sg:q:region}
For the radial oscillator evaluated at a general positive point,
or for the polynomial deformation in Part~\ref{bgg:part:arithmetic}, which
parameters give a complete sign pattern analogous to A321941?
\end{question}
The method of Sections~\ref{bgg:sg:sec:heat}--\ref{bgg:sg:sec:signproof}
suggests uniform coefficient majorants on
compact parameter sets.  The interesting targets are an explicit
parameter-dependent cutoff, the location of the exceptional
small indices, and a sharp boundary between sign regimes.
A separate algebraic problem is to relate that boundary to the
integer polynomial family already constructed in the repository.

\begin{question}[Complex continuation and exponentially small terms]
\label{bgg:sg:q:continuation}
What is the analytic continuation of \(\mathscr B_2(s)\) beyond
its disk of convergence, and how do its singularities determine
the exponentially small terms associated with \eqref{bgg:sg:eq:rho}?
\end{question}
Theorem~\ref{bgg:sg:thm:borel} proves a logarithmic singularity at
\(s=16\) from inside the disk.  It does not determine the full
singularity set, lateral continuations, or Stokes constants.
The exact heat integral is a natural starting point for such an
analysis.  Care is needed with the order-two factorial convention
and with the distinction between large-order coefficients and a
remainder theorem for an optimally truncated \(n\)-expansion.

\begin{question}[Arithmetic beyond the known congruence]
\label{bgg:sg:q:arithmetic}
Can the complete sign and large-order theorems be combined with
the repository's arithmetic recurrences to determine further
exact valuations and congruences of \(r_k\)?
\end{question}
The earlier report already proves integrality, the modulo-\(32\)
congruence, and sharp denominator conclusions in specified cases
(Part~\ref{bgg:part:arithmetic}).  Those are not new claims here.
The analytic scale can organize searches for additional arithmetic
patterns, but real asymptotics cannot by themselves determine a
2-adic valuation.  A proof would need a strengthened exact
recurrence or an appropriate congruence invariant.

\begin{question}[A formal proof with a small trusted certificate]
\label{bgg:sg:q:formal}
Which portions can be formalized first while keeping the
analytic dependency chain explicit?
\end{question}
The finite rational certificate, coefficientwise majorants, and
positive-kernel summation identity are good initial targets.
The exact spectral bridge requires a larger special-function and
operator-analysis library.  A useful staged formalization would
state the analytic inputs as explicit theorems, verify every
finite truncation algebraically, and then discharge the limiting
estimates.  This would also make the distinction between formal
series identities and convergent analytic identities mechanically
visible.

\emph{[Write note, batch 86A.]} The positive-kernel summation identity
named in the last question belongs to the manuscript's A088714 proof
(Part~V of that report); for A321941 the targets are the certificate of
Appendix~\ref{bgg:sg:app:certificate} and the majorants of
Section~\ref{bgg:sg:sec:majorants}.

% Part II's appendices continue Part I's letters.
\setcounter{section}{2}
\renewcommand{\thesection}{\Alph{section}}
\renewcommand{\theHsection}{sg.app.\Alph{section}}
\makeatletter\gdef\Hy@chapapp{\Hy@appendixstring}\makeatother

\section{A short exact certificate for the finite sign range}
\label{bgg:sg:app:certificate}

The following complete Python program proves the finite comparisons
used in Section~\ref{bgg:sg:sec:signproof}.  It constructs
\(B=q^{-1/2}\) from \(2qB'+q'B=0\), constructs
\(T=\tanh(t/4)\) from \(T'=(1-T^2)/4\), and constructs
\(E=e^{-T}\) from \(E'=-T'E\).
Truncated multiplication then implements \eqref{bgg:sg:eq:H}.
Every operation is exact.

\begin{lstlisting}[language=Python]
from fractions import Fraction as F
from math import factorial
N = 31

def mul(a, b):
    c = [F(0)] * (N + 1)
    for i, x in enumerate(a):
        for j, y in enumerate(b[:N + 1 - i]):
            c[i + j] += x * y
    return c

q = [F(0)] * (N + 1)
for j in range(N // 2 + 1):
    q[2*j] = F(1, 4**j * factorial(2*j + 1))
B = [F(1)]
for k in range(1, N + 1):
    B.append(-sum((F(k)-F(i, 2))*q[i]*B[k-i]
                  for i in range(1, k+1)) / k)
T = [F(0), F(1, 4)] + [F(0)] * (N - 1)
for k in range(2, N + 1):
    T[k] = -sum(T[i]*T[k-1-i] for i in range(k)) / (4*k)
E = [F(1)]
for k in range(1, N + 1):
    E.append(-sum(i*T[i]*E[k-i] for i in range(1, k+1))/k)
A = mul(B, E)
v = [F(0)] + q[:N]
power = [F(1)] + [F(0)] * N
S = power.copy()
b = F(3, 16)
for j in range(1, N + 1):
    if j > 1:
        b *= F((2*j-3)*(2*j+1), 16*j)
    power = mul(power, v)
    for k in range(j, N + 1):
        S[k] -= b * power[k]
h = mul(A, S)
assert h[:4] == [1, F(-7, 16), F(43, 1536), F(-61, 8192)]
assert all(h[k] < 0 for k in range(3, N + 1))
print("Exact certificate: h_k < 0 for 3 <= k <= 31.")
\end{lstlisting}

\emph{[Write note, batch 86A.]} This program is shipped as
\texttt{code/02-sign-finite\_certificate.py}, byte for byte the listing
above. It writes nothing; it prints one line.

\section{Source and dependency audit (Part~II)}\label{bgg:sg:app:sources}

The research reports were retrieved from the public repository at
commit
\begin{center}\small\ttfamily
b7e4f25b6e79cc669fb245cba520017a0f4228b7,
\end{center}
inspected on 4 October 2026.
The main published source is the final 2019 BGG article
\cite{BGG}; numbering in this manuscript refers to that final
version.

\subsection*{A321941 source}

The report \emph{Proving Two Arithmetic Conjectures for OEIS A321941}
is at
\begin{quote}\small
\nolinkurl{SetTheory/Cardinals/docs/reports/congruences-and-valuations/a321941-asymptotic-coefficient-integrality/article.tex}.
\end{quote}
Its Git blob identifier is
\begin{center}\small\ttfamily
10873a84fcf7a5f725ad74f90e2bd71ab7e7dd32.
\end{center}
Its proved integrality and congruence results are prior inputs to
the source comparison, but are not needed by the sign proof.
Its polynomial constant term equals \(-S_k\).
The contribution here is to prove the complete sign assertion,
identify that expression as the large-order scale, and derive
the all-orders and Borel conclusions.

\emph{[Write note, batch 86A.]} That report is
Part~\ref{bgg:part:arithmetic} of the present one; the blob is its
\texttt{article.tex} as it stood before this Part was added. The
manuscript's second source subsection, on the A088714 report, is printed
in Part~V of that report.

\subsection*{Verification boundaries}

The analytic and combinatorial arguments received separate
independent checks, including normalization constants, index
partitions, all-orders truncations, and the two infinite sums in
the convergence argument.  The executable checks and their
recorded outputs are included in the archive.
This internal review does not constitute external refereeing or
a proof-assistant certificate.  No OEIS entry or repository
commit is changed by this deliverable.

\emph{[Write note, batch 86A.]} ``Internal'' and ``independent'' refer to
reviews within the manuscript's own pipeline, recorded in its source
audit (shipped as
\path{data/05-bell-norm-SOURCE_AUDIT.json} of the A088714 report). The
two infinite sums of the convergence argument belong to its A088714
part.

\begin{thebibliography}{9}

\bibitem{BGG}
R. P. Brent, M. L. Glasser, and A. J. Guttmann,
\emph{A conjectured integer sequence arising from the exponential integral},
Journal of Integer Sequences \textbf{22} (2019), Article 19.4.7.
\href{https://cs.uwaterloo.ca/journals/JIS/VOL22/Brent/brent21.html}{Journal page};
\href{https://cs.uwaterloo.ca/journals/JIS/VOL22/Brent/brent21.pdf}{final journal PDF};
\href{https://arxiv.org/abs/1812.00316}{arXiv:1812.00316}.
Theorem numbering in this article follows the journal version.

\bibitem{oeis321941}
The OEIS Foundation Inc.,
\emph{A321941: Scaled numerators in the asymptotic expansion of the Maclaurin
coefficients in a Hadamard product involving the exponential integral}.
Entry contributed by Richard P. Brent, 2018.
\url{https://oeis.org/A321941}. Inspected 1 October 2026.

\bibitem{oeis262}
The OEIS Foundation Inc.,
\emph{A000262}, partitions into linearly ordered blocks (``sets of lists'').
\url{https://oeis.org/A000262}. Inspected 1 October 2026.

\bibitem{proveit-volume}
V. Reshetnikov and contributors, \emph{ProveIt},
\path{Analysis/Transseries/docs/series-and-transseries/Transseries_And_Inversion/README.md}.
Pinned commit: \repoSHA.
\href{https://github.com/VladimirReshetnikov/ProveIt/blob/0a6d798c5f39775b2fb4d85fef643fef996c87cf/Analysis/Transseries/docs/series-and-transseries/Transseries_And_Inversion/README.md}{Pinned repository source}.
Inspected 1 October 2026. The README, rather than the complete consolidated
volume, was used for the targeted comparison.

\bibitem{proveit-catalan}
V. Reshetnikov and contributors, \emph{ProveIt},
\path{Analysis/FabiusFunction/Lean/FabiusFunction/QuadraticCoreCatalan.lean},
same commit as \cite{proveit-volume}.
\href{https://github.com/VladimirReshetnikov/ProveIt/blob/0a6d798c5f39775b2fb4d85fef643fef996c87cf/Analysis/FabiusFunction/Lean/FabiusFunction/QuadraticCoreCatalan.lean}{Pinned source module}.
Source inspected; no new formal build performed for this article.

% Added for Part II (batch 86A); accessed by its manuscript on 4 October 2026.
\bibitem{dlmf-kummer}
NIST Digital Library of Mathematical Functions,
\emph{Section 13.2: Definitions and Basic Properties},
especially 13.2.34 and 13.2.40.
\url{https://dlmf.nist.gov/13.2}.
Accessed 4 October 2026.

\bibitem{dlmf-laguerre}
NIST Digital Library of Mathematical Functions,
\emph{Sections 18.3 and 18.18}, Laguerre orthogonality, completeness,
and the Hille--Hardy formula 18.18.27.
\url{https://dlmf.nist.gov/18.3};
\url{https://dlmf.nist.gov/18.18.E27}.
Accessed 4 October 2026.

\bibitem{dlmf-bessel}
NIST Digital Library of Mathematical Functions,
\emph{Sections 10.17 and 10.40}, asymptotic coefficients and
large-argument expansions of modified Bessel functions.
\url{https://dlmf.nist.gov/10.17.E1};
\url{https://dlmf.nist.gov/10.40.E1}.
Accessed 4 October 2026.

\bibitem{dlmf-watson}
NIST Digital Library of Mathematical Functions,
\emph{Section 2.3(ii): Watson's Lemma}.
\url{https://dlmf.nist.gov/2.3.ii}.
Accessed 4 October 2026.

\end{thebibliography}
\end{document}
